% =====================================================================
%  Rapport de stage — Typhoon : données climatiques par bâtiment
%  Stage d'été — SummerCamp 2026, Talan
%  Auteur : Hazem Mrad
%
%  Compilation : pdflatex main && pdflatex main
%  (aucune dependance BibTeX ; voir rapport/README.md pour les details)
% =====================================================================
\documentclass[11pt,a4paper,oneside]{report}

% ---------------------------------------------------------------------
%  Encodage, langue, police
%
%  Portabilite de moteur : pdfLaTeX garde le chemin historique
%  (inputenc + fontenc T1 + lmodern). Sous XeTeX (tectonic, XeLaTeX sur
%  Overleaf), inputenc est ignore par le moteur : les caracteres UTF-8
%  « - » et « œ » sont alors simplement ABSENTS de la police T1 et le
%  PDF sort ampute (299 glyphes manquants). fontspec est la solution
%  standard sous XeTeX. pdfLaTeX n'est pas modifie par ce garde-fou.
% ---------------------------------------------------------------------
\ifdefined\XeTeXversion
  \usepackage{fontspec}
\else
  \usepackage[utf8]{inputenc}
  \usepackage[T1]{fontenc}
  \usepackage{lmodern}
\fi
\usepackage[english,french]{babel}
\usepackage{csquotes}

% ---------------------------------------------------------------------
%  Mise en page
% ---------------------------------------------------------------------
\usepackage[a4paper,margin=2.4cm,top=2.6cm,bottom=2.6cm,headheight=15pt]{geometry}
\usepackage{fancyhdr}
\usepackage{titlesec}
\usepackage{caption}
\usepackage{subcaption}
\usepackage{enumitem}
\usepackage{float}
\usepackage{parskip}

% ---------------------------------------------------------------------
%  Mathématiques, tableaux, couleurs, graphiques
% ---------------------------------------------------------------------
\usepackage{amsmath,amssymb}
\usepackage{booktabs}
\usepackage{longtable}
\usepackage{tabularx}
\usepackage{array}
\usepackage{multirow}
\usepackage[table]{xcolor}
\usepackage{graphicx}
\usepackage{tikz}
\usetikzlibrary{arrows.meta,positioning,shapes.geometric,fit,backgrounds,calc}
\usepackage{pgfplots}
\pgfplotsset{compat=1.17}

% ---------------------------------------------------------------------
%  Code source
% ---------------------------------------------------------------------
\usepackage{listings}
\usepackage{textcomp}

% ---------------------------------------------------------------------
%  Encadrés
% ---------------------------------------------------------------------
\usepackage[most]{tcolorbox}

% =====================================================================
%  Palette (doit etre definie AVANT hyperref, qui l'utilise)
% =====================================================================
\definecolor{typhoonBlue}{HTML}{0B5FA5}
\definecolor{typhoonDeep}{HTML}{073E6C}
\definecolor{typhoonSand}{HTML}{F4F1EC}
\definecolor{typhoonGrey}{HTML}{6B7280}
\definecolor{typhoonGreen}{HTML}{1B7F5A}
\definecolor{typhoonAmber}{HTML}{B7791F}
\definecolor{typhoonRed}{HTML}{A6262C}
\definecolor{codeBg}{HTML}{F7F7F5}
\definecolor{codeKey}{HTML}{0B5FA5}
\definecolor{codeCom}{HTML}{6B7280}
\definecolor{codeStr}{HTML}{1B7F5A}

% ---------------------------------------------------------------------
%  Liens (à charger en dernier)
% ---------------------------------------------------------------------
\usepackage[hidelinks,colorlinks=true,linkcolor=typhoonBlue,
            citecolor=typhoonBlue,urlcolor=typhoonBlue]{hyperref}

% =====================================================================
%  Style des titres
% =====================================================================
\titleformat{\chapter}[display]
  {\normalfont\huge\bfseries\color{typhoonDeep}}
  {\filright\normalsize\sffamily\color{typhoonGrey}\MakeUppercase{\chaptertitlename}~\thechapter}
  {6pt}
  {\Huge\filright}
  [\vspace{4pt}{\color{typhoonBlue}\titlerule[1.2pt]}]

\titleformat{\section}
  {\normalfont\Large\bfseries\color{typhoonDeep}}{\thesection}{0.8em}{}
\titleformat{\subsection}
  {\normalfont\large\bfseries\color{typhoonBlue}}{\thesubsection}{0.8em}{}
\titleformat{\subsubsection}
  {\normalfont\normalsize\bfseries\color{typhoonGrey}}{\thesubsubsection}{0.8em}{}

\titlespacing*{\chapter}{0pt}{-10pt}{26pt}
\titlespacing*{\section}{0pt}{18pt}{8pt}
\titlespacing*{\subsection}{0pt}{13pt}{5pt}

% =====================================================================
%  En-têtes et pieds de page
% =====================================================================
\pagestyle{fancy}
\fancyhf{}
\fancyhead[L]{\small\sffamily\color{typhoonGrey}Typhoon — Rapport de stage}
\fancyhead[R]{\small\sffamily\color{typhoonGrey}\leftmark}
\fancyfoot[C]{\small\thepage}
\renewcommand{\headrulewidth}{0.4pt}
\renewcommand{\headrule}{\hbox to\headwidth{\color{typhoonBlue!40}\leaders\hrule height \headrulewidth\hfill}}
\fancypagestyle{plain}{\fancyhf{}\fancyfoot[C]{\small\thepage}\renewcommand{\headrulewidth}{0pt}}

% =====================================================================
%  Légendes
% =====================================================================
\captionsetup{
  font=small,
  labelfont={bf,color=typhoonDeep},
  labelsep=period,
  justification=raggedright,
  singlelinecheck=false
}

% =====================================================================
%  Listings
% =====================================================================
\lstdefinestyle{typhoon}{
  backgroundcolor=\color{codeBg},
  basicstyle=\ttfamily\footnotesize,
  keywordstyle=\color{codeKey}\bfseries,
  commentstyle=\color{codeCom}\itshape,
  stringstyle=\color{codeStr},
  numberstyle=\tiny\color{typhoonGrey},
  numbers=left,
  numbersep=8pt,
  frame=single,
  rulecolor=\color{typhoonBlue!25},
  framesep=6pt,
  breaklines=true,
  breakatwhitespace=false,
  postbreak=\mbox{\textcolor{typhoonBlue}{$\hookrightarrow$}\space},
  showstringspaces=false,
  tabsize=2,
  captionpos=b,
  xleftmargin=12pt,
  xrightmargin=6pt,
  columns=fullflexible,
  keepspaces=true,
  upquote=true
}
\lstset{style=typhoon}

\lstdefinelanguage{jsonc}{
  morecomment=[l]{//},
  morestring=[b]",
  literate=
    *{:}{{{\color{typhoonBlue}{:}}}}{1}
     {,}{{{\color{typhoonBlue}{,}}}}{1}
}
\lstdefinestyle{json}{
  style=typhoon,
  language=jsonc,
  numbers=none,
  morekeywords={true,false,null}
}

% =====================================================================
%  Encadrés
% =====================================================================
\newtcolorbox{invariant}[1][]{
  colback=typhoonBlue!5,
  colframe=typhoonBlue,
  boxrule=0.6pt,
  left=8pt,right=8pt,top=7pt,bottom=7pt,
  arc=2pt,
  fonttitle=\bfseries\sffamily,
  title=Invariant produit,
  #1
}
\newtcolorbox{honesty}[1][]{
  colback=typhoonAmber!7,
  colframe=typhoonAmber,
  boxrule=0.6pt,
  left=8pt,right=8pt,top=7pt,bottom=7pt,
  arc=2pt,
  fonttitle=\bfseries\sffamily,
  title=Honnêteté de la donnée,
  #1
}
\newtcolorbox{encadre}[1][]{
  colback=typhoonSand,
  colframe=typhoonGrey!60,
  boxrule=0.5pt,
  left=8pt,right=8pt,top=7pt,bottom=7pt,
  arc=2pt,
  #1
}

% =====================================================================
%  Raccourcis
% =====================================================================
\newcommand{\typhoon}{\textsc{Typhoon}}
\newcommand{\code}[1]{\texttt{\small #1}}
\newcommand{\fichier}[1]{\texttt{\small #1}}
\newcommand{\api}[1]{\texttt{\small #1}}
\newcommand{\frref}[1]{FR-#1}
\newcommand{\scn}[1]{SCN-#1}
\newcommand{\vfx}[1]{VFX-#1}

% =====================================================================
%  MÉTADONNÉES DU RAPPORT
%  ---------------------------------------------------------------------
%  Toutes les informations de la page de garde sont regroupées ici.
%  Modifier UNIQUEMENT ce bloc pour mettre le rapport à jour.
% =====================================================================

% --- Auteur -----------------------------------------------------------
\newcommand{\reportAuthor}{Hazem Mrad}

% --- Établissement ----------------------------------------------------
\newcommand{\reportSchool}{ESPRIT}
\newcommand{\reportSchoolLong}{École Supérieure Privée d'Ingénierie et de Technologie}

% --- Intitulé du diplôme (à adapter à votre cursus) --------------------
\newcommand{\reportDegree}{Diplôme d'ingénieur}
\newcommand{\reportCycle}{Cycle d'ingénieur}

% --- Entreprise d'accueil ---------------------------------------------
\newcommand{\reportCompany}{Talan}
\newcommand{\reportCompanyLong}{Groupe Talan --- Conseil et ingénierie technologique}
\newcommand{\reportProgramme}{SummerCamp 2026}

% --- Période de stage --------------------------------------------------
\newcommand{\reportPeriod}{14 août 2026 -- 15 septembre 2026}

% --- Année universitaire ----------------------------------------------
\newcommand{\reportAcademicYear}{2026--2027}

% --- Tuteurs ----------------------------------------------------------
%  Laisser les pointillés, ou remplacer par les noms.
\newcommand{\reportTutorAcademic}{\dotfill}
\newcommand{\reportTutorCompany}{\dotfill}

% ---------------------------------------------------------------------
%  LOGO DE L'ÉCOLE
%  Déposez le fichier sous rapport/logos/esprit.png (ou .jpg / .pdf).
%  Tant qu'aucun fichier n'est présent, un cadre de remplacement
%  explicite est affiché : le rapport compile quand même.
% ---------------------------------------------------------------------
\newcommand{\reportSchoolLogo}[1][2.4cm]{%
  \IfFileExists{logos/esprit.png}{\includegraphics[height=#1]{logos/esprit.png}}{%
  \IfFileExists{logos/esprit.jpg}{\includegraphics[height=#1]{logos/esprit.jpg}}{%
  \IfFileExists{logos/esprit.pdf}{\includegraphics[height=#1]{logos/esprit.pdf}}{%
    \begin{tikzpicture}
      \node[draw=typhoonGrey,dashed,rounded corners=3pt,inner sep=9pt,
            text width=2.4cm,align=center,font=\sffamily\scriptsize,
            text=typhoonGrey,minimum height=#1]
        {Logo de l'école\\[2pt]à déposer dans\\[2pt]\texttt{logos/}};
    \end{tikzpicture}}}}%
}

% ---------------------------------------------------------------------
%  LOGO TALAN (vectorisé depuis le SVG officiel du dépôt :
%  frontend/public/vetter-consulting/assets/img/talan-logo-white.svg)
%  Tracé blanc, à poser sur un fond sombre.
% ---------------------------------------------------------------------
\newcommand{\reportCompanyLogo}{%
\begin{tikzpicture}[x=0.0032cm,y=-0.0032cm]
  \fill[white]
    (1160,0) -- (1160,280.78) -- (1084.14,280.78) -- (968.48,154.207) --
    (968.48,280.78) -- (905.856,280.78) -- (875.875,203.506) -- (875.875,0) --
    (951.732,0) -- (1067.39,132.409) -- (1067.39,0) -- cycle;
  \fill[white]
    (875.875,203.518) -- (905.856,280.791) -- (875.875,280.791) -- cycle;
  \fill[white]
    (875.887,203.506) -- (875.887,280.767) -- (806.954,280.767) --
    (795.636,246.812) -- (692.132,246.812) -- (680.394,280.767) --
    (640.931,280.767) -- (667.791,202.811) -- (611.75,202.811) --
    (690.467,0) -- (796.908,0) -- cycle
    (770.088,171.387) -- (744.107,95.1234) -- (718.126,171.387) -- cycle;
  \fill[white]
    (667.784,202.828) -- (640.924,280.784) -- (581.5,280.784) --
    (611.743,202.828) -- cycle;
  \fill[white]
    (611.744,202.824) -- (581.501,280.78) -- (523.978,280.78) --
    (493.656,202.614) -- (493.656,0) -- (586.275,0) -- (586.275,202.824) --
    cycle;
  \fill[white]
    (493.656,202.619) -- (523.978,280.785) -- (493.656,280.785) -- cycle;
  \fill[white]
    (493.646,202.614) -- (493.646,280.767) -- (425.08,280.767) --
    (413.762,246.812) -- (310.258,246.812) -- (298.52,280.767) --
    (199.633,280.767) -- (278.336,77.9427) -- (308.579,77.9427) --
    (308.579,0) -- (415.034,0) -- (493.659,202.614) -- cycle
    (388.188,171.387) -- (362.22,95.1234) -- (336.239,171.387) -- cycle;
  \fill[white]
    (308.575,0) -- (308.575,77.9558) -- (278.332,77.9558) -- cycle;
  \fill[white]
    (308.571,0) -- (278.327,77.9558) -- (199.624,77.9558) --
    (199.624,280.78) -- (107.006,280.78) -- (107.006,77.9558) --
    (0,77.9558) -- (0,0) -- cycle;
\end{tikzpicture}%
}

% =====================================================================
%  Métadonnées PDF
% =====================================================================
\hypersetup{
  pdftitle={Typhoon — Service de données climatiques par bâtiment},
  pdfauthor={Hazem Mrad},
  pdfsubject={Rapport de stage — SummerCamp 2026, Talan},
  pdfkeywords={assurance, climat, Géorisques, BDNB, FastAPI, React, SIG}
}

% =====================================================================
%  DOCUMENT
% =====================================================================
\begin{document}

\selectlanguage{french}

% ---------------------------------------------------------------------
%  Page de garde — conforme aux usages académiques français
% ---------------------------------------------------------------------
\begin{titlepage}
\thispagestyle{empty}

% Bandeau supérieur (logo Talan en blanc sur fond sombre)
\begin{tikzpicture}[remember picture,overlay]
  \fill[typhoonDeep] (current page.north west) rectangle
      ([yshift=-3.5cm]current page.north east);
  \fill[typhoonBlue] ([yshift=-3.5cm]current page.north west) rectangle
      ([yshift=-3.56cm]current page.north east);
\end{tikzpicture}

\noindent\begin{minipage}[c]{0.46\textwidth}
  \color{white}\sffamily
  {\bfseries\large \reportCompany}\\[2pt]
  {\footnotesize \reportProgramme}\\[2pt]
  {\scriptsize\color{white!75}\reportCompanyLong}
\end{minipage}\hfill
\begin{minipage}[c]{0.30\textwidth}
  \raggedleft\reportCompanyLogo
\end{minipage}

\vspace{1.5cm}

% --- Établissement ---------------------------------------------------
\begin{center}
  \reportSchoolLogo[2.3cm]\\[10pt]
  {\sffamily\bfseries\large \reportSchool}\\[2pt]
  {\sffamily\small\color{typhoonGrey} \reportSchoolLong}
\end{center}

\vspace{1.2cm}

% --- Nature et titre du travail --------------------------------------
\begin{center}
  {\sffamily\small\color{typhoonGrey}\MakeUppercase{Rapport de stage de fin de cycle}}\\[6pt]
  {\color{typhoonBlue}\rule{4cm}{0.8pt}}\\[14pt]

  {\sffamily\bfseries\LARGE\color{typhoonDeep}
   Service de données climatiques\\[4pt] par bâtiment}\\[12pt]

  {\sffamily\large\color{typhoonDeep}
   Concevoir et industrialiser une jointure aléa $\times$ vulnérabilité\\
   à l'échelle du bâtiment pour les assureurs français}\\[10pt]

  {\sffamily\small\color{typhoonGrey}
   De la maquette de démonstration à un contrat de données traçable,\\
   testé et déployé sur une infrastructure serverless.}
\end{center}

\vfill

% --- Identification ---------------------------------------------------
\begin{center}
\renewcommand{\arraystretch}{1.45}
\begin{tabular}{@{}p{4.6cm}@{\hspace{4pt}}p{7.6cm}@{}}
  \textcolor{typhoonGrey}{\sffamily\small Présenté par} &
    \textbf{\reportAuthor} \\
  \textcolor{typhoonGrey}{\sffamily\small Intitulé du diplôme} &
    \reportDegree \\
  \textcolor{typhoonGrey}{\sffamily\small Cycle} &
    \reportCycle \\
  \textcolor{typhoonGrey}{\sffamily\small Structure d'accueil} &
    \reportCompanyLong \\
  \textcolor{typhoonGrey}{\sffamily\small Période de stage} &
    \reportPeriod \\
  \textcolor{typhoonGrey}{\sffamily\small Tuteur pédagogique} &
    \reportTutorAcademic \\
  \textcolor{typhoonGrey}{\sffamily\small Tuteur entreprise} &
    \reportTutorCompany \\
  \textcolor{typhoonGrey}{\sffamily\small Année universitaire} &
    \textbf{\reportAcademicYear} \\
\end{tabular}
\end{center}

\vfill

\begin{center}
  {\color{typhoonBlue}\rule{\textwidth}{0.6pt}}\\[6pt]
  {\sffamily\footnotesize\color{typhoonGrey}
   Rapport de stage --- \reportAcademicYear \hspace{2em}·\hspace{2em}
   \reportAuthor \hspace{2em}·\hspace{2em} \reportCompany}
\end{center}

\end{titlepage}

% ---------------------------------------------------------------------
%  Remerciements
% ---------------------------------------------------------------------
\chapter*{Remerciements}
\addcontentsline{toc}{chapter}{Remerciements}
\markboth{Remerciements}{Remerciements}

Ce stage d'été n'aurait pas pris cette forme sans plusieurs personnes que je
souhaite remercier ici.

Je remercie \textbf{Talan} de m'avoir accueilli dans le cadre du SummerCamp
2026 et de m'avoir laissé une latitude réelle sur le sujet : proposer une
problématique, la reformuler, puis la construire jusqu'au déploiement. Un stage
où l'on peut supprimer du code parce qu'il ne sert pas le produit vaut plus
qu'un stage où l'on empile des fonctionnalités.

Je remercie les personnes de l'écosystème assurance — l'actuaire et le
professionnel de l'assurance qui ont accepté de réagir au projet au démarrage.
Leurs réponses, parfois dérangeantes, ont directement réécrit la direction
produit : le \og score \fg{} que je pensais être le livrable a été explicitement
refusé, et cette conversation a probablement été le meilleur investissement des
deux mois.

Je remercie enfin les équipes et les projets open data qui rendent ce travail
possible : Géorisques, la BDNB du CSTB, la Géoplateforme de l'IGN, Sandre,
Panoramax, Open-Meteo et le Copernicus Climate Data Store. Un service comme
\typhoon{} ne remplace aucune de ces sources : il existe parce qu'elles sont
publiques, documentées et réutilisables.

\vspace{1em}
\begin{flushright}
\textit{Hazem Mrad}
\end{flushright}

% ---------------------------------------------------------------------
%  Résumé
% ---------------------------------------------------------------------
\chapter*{Résumé}
\addcontentsline{toc}{chapter}{Résumé}
\markboth{Résumé}{Résumé}

Ce rapport présente le travail réalisé pendant un stage d'été au sein du
programme SummerCamp 2026 de Talan. Le sujet est un service de données
climatiques \textbf{par bâtiment} destiné aux assureurs français, baptisé
\typhoon{}.

Le problème de départ est concret. Les données publiques existent et sont
gratuites : Géorisques publie les aléas réglementaires, la BDNB décrit les
bâtiments, Copernicus produit des projections climatiques. Mais aucune de ces
sources ne relie l'exposition d'un bâtiment à sa vulnérabilité. Géorisques ne
sait pas ce qu'est le bâtiment ; la BDNB ne sait pas s'il est dans un périmètre
d'inondation. Aujourd'hui, un souscripteur résout ce manque en ouvrant deux
onglets et en jugeant à l'œil.

La contribution du stage tient en trois apports. D'abord, une
\textbf{reformulation produit} : le livrable n'est pas une interface, c'est un
contrat de données unique, l'adresse en entrée, un JSON canonique en sortie, où
chaque fait porte sa source et son niveau de résolution, et où aucun score
composite n'apparaît — la notation reste le métier de l'assureur. Ensuite, une
\textbf{industrialisation de la jointure} : un connecteur WFS Géorisques
(point-in-polygon écrit à la main, parseur GML 3.2) croisé avec un passage
verbatim de la BDNB, 13 aléas normalisés, un cache, un quota, une
authentification par clé et une suite de tests hors-ligne. Enfin, une
\textbf{couche de mise en situation} sur la carte : des classes de retour
officielles du TRI, une timeline horaire, un moteur où le scénario pilote le
niveau et la pluie n'en pilote que la forme, et un mode de rendu cinématique
explicitement étiqueté comme non contractuel.

Le résultat est une application déployée et vérifiée en ligne, un backend de
70 fichiers Python, un frontend de 68 modules TypeScript, \textbf{214 tests
backend et 143 tests frontend} hors-ligne, une constitution d'ingénierie qui
interdit les scores dans le contrat sérialisé, et un ensemble de décisions
documentées — y compris celles qui ont consisté à \emph{supprimer} du code
livré.

\vspace{1em}
\noindent\textbf{Mots-clés :} assurance, risque climatique, open data,
Géorisques, BDNB, TRI, jointure spatiale, FastAPI, React, Mapbox GL JS,
serverless.

\clearpage

% =====================================================================
%  Synthèse — fiche en une page
% =====================================================================
\chapter*{Synthèse en une page}
\addcontentsline{toc}{chapter}{Synthèse en une page}
\markboth{Synthèse}{Synthèse}

\begin{center}
\begin{tcolorbox}[colback=typhoonDeep,colframe=typhoonBlue!0,arc=3pt,
  left=10pt,right=10pt,top=8pt,bottom=8pt,width=\textwidth]
\centering
{\color{white}\sffamily\bfseries\large \typhoon{} — données climatiques par
bâtiment}\\[3pt]
{\color{white!85}\sffamily\small Une adresse, treize aléas, un bâtiment, une
provenance sur chaque fait.}
\end{tcolorbox}
\end{center}

\section*{Le problème}

Géorisques publie les périmètres d'aléas mais ne connaît pas le bâtiment. La
BDNB décrit les bâtiments mais ne sait pas lesquels sont inondables. Aucune
source ne relie exposition et vulnérabilité à l'échelle du bâtiment. Un
souscripteur comble ce vide à l'œil, sans traçabilité.

\section*{La réponse}

Un service qui prend une adresse française et renvoie un contrat JSON
canonique : treize aléas normalisés, le passage verbatim de la BDNB, une
provenance et un niveau de résolution sur chaque fait — et aucun score.

\section*{Les trois décisions structurantes}

\begin{enumerate}[leftmargin=1.5em,itemsep=3pt]
  \item \textbf{Refuser le score composite.} Sous Solvabilité~II, un score
        tiers engage son producteur et se fait reconstruire ou refuser par
        l'actuaire. Le champ de résolution devient le signal de qualité.
  \item \textbf{Faire de l'incertitude une fonctionnalité.} Trois états
        (présent / absent / inconnu) et une distinction stricte entre absence
        de risque et panne de service.
  \item \textbf{Soustraire plutôt qu'ajouter.} La phase la plus productive du
        stage a consisté à supprimer quatre connecteurs, une couche de
        notation et un tableau de bord.
\end{enumerate}

\section*{Chiffres clés}

\begin{center}
\begin{tabular}{@{}r@{\hspace{6pt}}l@{\hspace{18pt}}r@{\hspace{6pt}}l@{}}
\textbf{357} & tests automatisés hors-ligne & \textbf{28} & exigences fonctionnelles \\
\textbf{214} & tests backend (pytest) & \textbf{18} & tâches planifiées (+10 scénarios) \\
\textbf{143} & tests frontend (Vitest) & \textbf{13} & aléas normalisés \\
\textbf{70} & fichiers Python (backend) & \textbf{4} & aléas certifiés \\
\textbf{68} & modules TypeScript (frontend) & \textbf{9} & sources externes intégrées \\
\end{tabular}
\end{center}

\section*{Résultat}

Un service déployé et vérifié en ligne le 15 septembre 2026, couvrant
l'ensemble du parcours : géocodage, jointure géométrique, fiche de
vulnérabilité, mise en situation cartographique et export. Le prototype
démontre une thèse : la donnée publique française permet de produire une
information de risque climatique par bâtiment, traçable et défendable — à
condition d'accepter de ne pas tout savoir et de le dire.

\vfill

\begin{tcolorbox}[colback=typhoonAmber!8,colframe=typhoonAmber,arc=2pt,
  left=8pt,right=8pt,top=6pt,bottom=6pt]
{\sffamily\bfseries\small Ce que le prototype n'est pas :}
{\small un produit commercialisable. L'authentification est simulée, le
magasin de lots est en mémoire, la certification attend sa qualification
humaine et le budget de latence n'est pas calibré. Ces limites sont détaillées
au chapitre~9 et assumées.}
\end{tcolorbox}

\clearpage

% ---------------------------------------------------------------------
%  Sommaire
% ---------------------------------------------------------------------
\tableofcontents
\clearpage
\listoffigures
\clearpage
\listoftables
\clearpage

% ---------------------------------------------------------------------
%  Chapitres
% ---------------------------------------------------------------------
% =====================================================================
\chapter{Contexte du stage}
\label{ch:contexte}
% =====================================================================

\section{La structure d'accueil}

Talan est un groupe de conseil et d'ingénierie technologique dont l'activité
couvre le conseil en organisation, l'intégration de systèmes et, de plus en
plus, la \emph{data} et l'intelligence artificielle. Ses interventions
s'adressent à plusieurs secteurs, dont la banque, l'assurance, l'énergie et les
transports.

Ce stage s'est déroulé dans le cadre du \textbf{SummerCamp 2026}, le programme
d'été du groupe. Le format est celui d'un stage projet : un sujet ouvert, une
durée courte — environ un mois de travail effectif, du 14 août au 15 septembre
2026 d'après l'historique du dépôt — et une obligation de résultat tangible.
Contrairement à un stage d'observation, il fallait livrer quelque chose qui
tourne : un service déployé, pas une présentation.

Le sujet proposé portait sur le \textbf{risque climatique appliqué à
l'assurance}, un thème qui n'est pas un exercice académique. La France a
effectivement ouvert une discussion de politique publique sur l'adaptation du
régime d'indemnisation des catastrophes naturelles (CatNat) au changement
climatique. Travailler sur la donnée climatique appliquée au bâti, dans ce
contexte, n'est pas un problème inventé pour un stage : c'est un problème que
des professionnels cherchent activement à outiller.

\section{Le problème posé}

\subsection{Trois sources, aucune jointure}

Le point de départ du projet est un constat simple, mais qui n'avait pas été
exploité. La France dispose d'un écosystème open data remarquable sur le
risque :

\begin{itemize}[leftmargin=1.4em]
  \item \textbf{Géorisques} publie les aléas réglementaires : périmètres de
        plans de prévention des risques (PPR), sites et sols pollués, réseaux
        de transport de matières dangereuses, cavités, retrait-gonflement des
        argiles, etc. La donnée est servie en WFS, donc interrogeable
        géométriquement.
  \item La \textbf{BDNB} (Base de Données Nationale des Bâtiments, CSTB) décrit
        chaque bâtiment : usage, hauteur, emprise, matériaux de construction,
        année, performance énergétique, avec des indicateurs de fiabilité.
  \item \textbf{Copernicus} et \textbf{Open-Meteo} fournissent la dimension
        temporelle : projections climatiques à 2050 et 2100, prévisions
        météorologiques horaires, débits de rivière.
\end{itemize}

Aucune de ces sources ne fait la jointure. Géorisques ne sait pas ce qu'est le
bâtiment : il dit qu'un périmètre existe, pas si \emph{ce} bâtiment est dedans.
La BDNB ne sait pas si le bâtiment est inondable. Copernicus travaille sur une
grille d'environ 10~km, donc ne dit rien d'un bâtiment précis.

\begin{encadre}
\textbf{La jointure aléa $\times$ vulnérabilité est le produit.} Un aléa sans
vulnérabilité est une carte. Une vulnérabilité sans aléa est une fiche
technique. Les deux ensemble, pour une même adresse, avec la provenance de
chaque fait, constituent une information qui n'existe nulle part ailleurs.
\end{encadre}

\subsection{Le contournement actuel}

Un souscripteur qui reçoit un dossier d'assurance habitation ou professionnelle
doit aujourd'hui évaluer l'exposition du bien. La pratique observée — décrite
dans les notes de cadrage du projet et confirmée lors des échanges avec un
professionnel du secteur — consiste à ouvrir Géorisques dans un onglet, la BDNB
dans un autre, et à juger visuellement : la rivière est-elle proche ? le
périmètre enveloppe-t-il la parcelle ? le bâtiment est-il en rez-de-chaussée ?

C'est lent, non reproductible, non traçable, et cela ne passe pas l'épreuve
d'un audit. Sous Solvabilité~II, une décision de souscription doit pouvoir être
justifiée. \og J'ai regardé la carte et ça m'a paru bon \fg{} n'est pas une
justification.

\section{La mission}

La mission confiée peut se formuler ainsi :

\begin{quote}
\itshape
Concevoir un service qui, à partir d'une adresse française, produit une fiche
de risque climatique \textbf{au niveau du bâtiment}, en croisant les aléas
réglementaires et la vulnérabilité du bâti, avec une provenance explicite pour
chaque fait — et le rendre utilisable par un souscripteur dans son travail
quotidien.
\end{quote}

Cette formulation est volontairement large. Elle laisse ouvertes trois
questions qui se sont révélées être les vraies décisions du stage :

\begin{enumerate}[leftmargin=1.6em]
  \item \textbf{Quel est le livrable ?} Une interface graphique, un rapport
        narratif, ou une donnée structurée ? La réponse n'était pas évidente
        au départ.
  \item \textbf{Quel est le niveau de résolution acceptable ?} Faire du
        per-bâtiment partout est-il possible ? Sinon, comment le dire
        honnêtement ?
  \item \textbf{Que fait-on de la synthèse ?} Un assureur attend souvent un
        \og score \fg. Est-ce le bon livrable ?
\end{enumerate}

\section{Démarche et organisation du rapport}

Le stage s'est déroulé en trois mouvements, qui structurent ce rapport.

\begin{description}[leftmargin=0em,style=nextline]
  \item[Phase 1 — Exploration (mi-août 2026).] Construction d'une première
        version fonctionnelle : quatre connecteurs de données, une couche de
        notation, un tableau de bord climatique, des extrusions 3D. Le produit
        \emph{marchait}, mais son intérêt était difficile à défendre : il
        reproduisait, à résolution égale et gratuitement, ce que chaque portail
        source fournissait déjà.
  \item[Phase 2 — Recentrage (25 août 2026).] Une remise à plat complète. Le
        livrable a été redéfini comme un contrat de données, le score composite
        a été supprimé, et une grande partie du code de la phase 1 a été
        \emph{supprimée}. Une constitution d'ingénierie a été écrite pour
        verrouiller ces choix par des tests et des règles non négociables.
  \item[Phase 3 — Industrialisation et mise en scène (fin août à
        mi-septembre).] Durcissement du contrat, ajout d'un moteur de scénarios
        interactifs sur la carte, d'un mode de rendu cinématique, et d'un
        déploiement serverless vérifié en ligne.
\end{description}

Le rapport suit cette logique. Le chapitre~\ref{ch:metier} décrit le métier de
l'assurance et ce que le produit doit — ou ne doit pas — être. Le
chapitre~\ref{ch:existant} analyse l'existant et la décision de recentrage. Les
chapitres~\ref{ch:methodo} à \ref{ch:architecture} traitent de la méthode, des
sources de données et de l'architecture. Les chapitres~\ref{ch:scenarios} et
\ref{ch:vfx} décrivent le moteur de scénarios et le mode VFX. Le
chapitre~\ref{ch:qualite} détaille la stratégie de test. Le
chapitre~\ref{ch:deploiement} couvre le déploiement, le
chapitre~\ref{ch:resultats} les résultats mesurés, et le
chapitre~\ref{ch:bilan} le bilan critique.

\begin{figure}[H]
\centering
\begin{tikzpicture}[
  node distance=8mm,
  phase/.style={draw=typhoonBlue,thick,fill=typhoonBlue!6,rounded corners=3pt,
                text width=3.55cm,align=center,minimum height=2.5cm,inner sep=6pt},
  arrow/.style={-{Latex[length=2.4mm]},thick,typhoonBlue}
]
\node[phase] (p1) {\textbf{\color{typhoonDeep}Phase 1}\\[3pt]
  {\small Exploration}\\[4pt]
  {\footnotesize 4 connecteurs, notation, tableau de bord, 3D}\\[3pt]
  {\scriptsize\color{typhoonGrey}14--24 août}};
\node[phase,right=of p1] (p2) {\textbf{\color{typhoonDeep}Phase 2}\\[3pt]
  {\small Recentrage}\\[4pt]
  {\footnotesize Contrat de données, suppression du score, constitution}\\[3pt]
  {\scriptsize\color{typhoonGrey}25--31 août}};
\node[phase,right=of p2] (p3) {\textbf{\color{typhoonDeep}Phase 3}\\[3pt]
  {\small Industrialisation}\\[4pt]
  {\footnotesize Scénarios TRI, mode VFX, tests, déploiement Vercel}\\[3pt]
  {\scriptsize\color{typhoonGrey}1--15 sept.}};
\draw[arrow] (p1) -- (p2);
\draw[arrow] (p2) -- (p3);
\end{tikzpicture}
\caption{Les trois phases du stage. La phase 2 est une phase de soustraction :
elle retire davantage qu'elle n'ajoute, et c'est ce qui a rendu le produit
défendable.}
\label{fig:phases}
\end{figure}

% =====================================================================
\chapter{Le métier de l'assurance et les invariants du produit}
\label{ch:metier}
% =====================================================================

Un projet technique appliqué à un métier se juge d'abord sur sa compréhension
de ce métier. Ce chapitre expose ce que j'ai appris du secteur, comment cela a
déterminé les choix de conception, et les règles qui en découlent — règles que
le projet s'interdit de violer.

\section{Deux acheteurs, un moteur}

Les échanges menés au démarrage du stage ont fait apparaître que \og
l'assureur \fg{} n'est pas un utilisateur unique. Deux profils aux besoins
différents cohabitent.

\subsection{Le souscripteur (\emph{underwriter})}

C'est l'utilisateur principal de la version~1. Il traite des dossiers : une
adresse arrive, il doit décider \emph{accepter}, \emph{mettre à l'étude} ou
\emph{refuser}. Ses contraintes :

\begin{itemize}[leftmargin=1.4em]
  \item \textbf{rapidité} — la décision s'insère dans un flux de travail, pas
        dans une analyse hebdomadaire ;
  \item \textbf{jugement spatial} — il pense en distance, en périmètre, en
        relief. Il regarde une carte, il ne lit pas un tableau de bord ;
  \item \textbf{traçabilité} — sa décision peut être auditée. Il doit pouvoir
        dire d'où vient chaque élément.
\end{itemize}

C'est pour ce profil que la carte reste une fonctionnalité de premier plan,
et non une décoration : elle ne montre pas les couches sources, elle montre
\emph{la jointure} — l'empreinte du bâtiment colorée par le résultat du test
géométrique, accompagnée d'une fiche de vulnérabilité.

\subsection{L'actuaire (\emph{pricing})}

Profil secondaire, hors périmètre de la version~1 mais structurant. Il
construit des modèles de tarification et de provisionnement. Ses besoins sont
presque opposés :

\begin{itemize}[leftmargin=1.4em]
  \item \textbf{pas d'interface} — il veut de la donnée, pas un écran ;
  \item \textbf{volume} — il traite des portefeuilles entiers, pas une adresse ;
  \item \textbf{variables brutes} — il veut appliquer \emph{ses} formules.
\end{itemize}

Le point décisif : l'actuaire interrogé ne voulait \emph{pas} de notre score.
Il voulait les variables par aléa, par horizon et par scénario climatique,
étiquetées et propres, pour exécuter ses propres calculs. Cette réponse a
directement conduit à l'invariant central du produit.

\section{La non-stationnarité : pourquoi la donnée historique ne suffit pas}

Le raisonnement qui fonde l'intérêt du produit mérite d'être explicité, car il
n'est pas immédiat.

Un assureur possède un actif que nous n'aurons jamais : l'historique de ses
sinistres. Pour un bien déjà couvert, cet historique est un meilleur prédicteur
que n'importe quelle reconstitution à partir de sources publiques. Il serait
donc absurde de prétendre \og mieux connaître le risque actuel \fg.

Mais le risque climatique est \textbf{non stationnaire} : le futur ne ressemble
pas au passé. Un historique de sinistres, aussi volumineux soit-il, décrit des
pertes réalisées sous le climat d'hier. Il ne peut structurellement pas
indiquer comment l'exposition d'un bâtiment donné évolue sur les 25 à 75
prochaines années. Cette information doit venir de la science du climat et des
modèles de projection, pas d'un registre de sinistres.

\begin{encadre}
\textbf{L'argument de vente n'est donc pas \og plus de données \fg.} C'est \og
la partie de l'information qui ne peut pas venir de votre historique \fg.
Aucune quantité de données propriétaires ne comble ce trou précis.
\end{encadre}

C'est pour cette raison que la trajectoire temporelle — aujourd'hui, 2050,
2100 — est le cœur du produit, et non un ornement. Et c'est pour cette raison
que la donnée climatique, un temps écartée du périmètre, a dû y revenir sous
une forme honnête (voir chapitre~\ref{ch:donnees}).

\section{Les invariants du produit}

Le projet a produit un document de référence, la \emph{constitution}
(\fichier{constitution.md}), qui fixe les règles non négociables. Chaque règle
porte un mécanisme d'application : un test automatisé, une porte CI, ou une
revue humaine. En voici les principales, avec leur justification.

\subsection{Faits et provenance uniquement — pas de score}

\begin{invariant}
Aucune réponse de l'API ne peut contenir de score composite, de jugement de
risque, de recommandation, ni d'évaluation générée par IA. La graduation
fournie par la source est transmise verbatim sous forme de texte ; toute
conversion en bandes est de la coloration de présentation et n'est
\emph{jamais} sérialisée dans le contrat canonique.
\end{invariant}

La justification est double. Juridiquement, un score tiers de risque
climatique, produit par un acteur qui n'est ni l'assureur ni le régulateur,
crée une responsabilité : s'il est erroné et qu'il a servi à refuser un
contrat, qui est responsable ? Techniquement, un score se fait
\emph{reconstruire} ou \emph{refuser} : l'actuaire qui reçoit un score
non transparent finit par ignorer la boîte noire ou par la réimplémenter.

La conséquence architecturale est forte : puisque le score disparaît, le champ
\texttt{resolution} devient \emph{le seul signal de qualité} du contrat. Chaque
objet aléa doit dire à quel niveau il a été établi.

\subsection{Un niveau de résolution sur chaque fait}

\begin{invariant}
Chaque objet aléa du contrat porte \texttt{source}, \texttt{resolution} et les
champs d'attribution exigés par la Licence Ouverte~2.0. Un objet aléa sans ces
champs échoue la validation du schéma.
\end{invariant}

Trois valeurs de résolution existent, et elles ne sont pas interchangeables :

\begin{table}[H]
\centering
\small
\caption{Les trois niveaux de résolution et leur signification exacte.}
\label{tab:resolution}
\begin{tabularx}{\textwidth}{@{}l l X@{}}
\toprule
\textbf{Valeur} & \textbf{Origine} & \textbf{Signification} \\
\midrule
\code{per-building} & Test géométrique WFS &
  Le point a été testé contre un polygone source. L'affirmation porte
  réellement sur ce bâtiment. \\
\addlinespace
\code{commune-level} & Décret ou arrêté &
  Le fait est communal \emph{par nature} (zonage sismique, classement radon).
  Ce n'est pas une dégradation : c'est la résolution maximale qui existe. \\
\addlinespace
\code{commune-level-estimate} & Repli, ou RGA &
  Nous n'avons pas pu établir le fait au niveau du bâtiment. Le fait est une
  estimation communale. À traiter avec prudence. \\
\bottomrule
\end{tabularx}
\end{table}

La distinction entre les deux derniers cas est essentielle. Un aléa sismique
\code{commune-level} est \emph{correct} : le zonage est décrétal, il n'existe
pas de version par bâtiment. Un aléa inondation \code{commune-level-estimate}
signale au contraire que la source n'a pas répondu et que l'information est
dégradée. Confondre les deux ferait passer une panne pour un fait.

\subsection{Passage verbatim de la BDNB}

\begin{invariant}
Le contrat transmet les attributs BDNB sans suppression ni renommage
silencieux. Un champ indisponible est explicitement \texttt{null}, avec
l'échec enregistré dans la liste d'erreurs — jamais omis.
\end{invariant}

C'est une conséquence directe du profil actuaire : s'il doit filtrer lui-même,
il lui faut l'intégralité du champ. La curation — quels attributs montrer dans
le panneau latéral — se fait uniquement à la couche de présentation.

\subsection{Une seule forme canonique}

\begin{invariant}
Il existe exactement un schéma de réponse par ressource. Deux formes vivantes
pour la même ressource constituent un défaut.
\end{invariant}

Cette règle a une histoire : au cours du projet, deux pipelines concurrents
produisaient deux JSON différents pour la même adresse. L'un était utilisé par
l'interface, l'autre par le lot. Les divergences silencieuses entre les deux
sont exactement le genre de bug qui détruit la confiance dans un contrat de
données. La règle a été formalisée en question ouverte (OQ-1) puis tranchée.

\subsection{Le lot est N fois l'unitaire}

\begin{invariant}
\api{POST /diagnostic/batch} traite chaque adresse à travers le pipeline
unitaire identique. Aucune seconde implémentation en masse, aucun chemin de
code d'extraction versionnée.
\end{invariant}

L'intérêt est la garantie de cohérence : il ne peut pas exister de différence
entre le résultat d'un appel unitaire et celui du même appel au sein d'un lot,
hors horodatage. Cette équivalence est testée.

\subsection{Honnêteté sur le RGA}

\begin{invariant}
Tant qu'aucune source vectorielle n'est câblée, le retrait-gonflement des
argiles (RGA) est livré avec \code{resolution: "commune-level-estimate"} et
rien ne peut revendiquer une résolution par bâtiment pour cet aléa.
\end{invariant}

Le RGA est un motif majeur de sinistres en France. La tentation d'afficher un
résultat par bâtiment est donc forte. Mais la donnée disponible au moment du
stage était communale. Le choix a été de l'annoncer plutôt que de la surinterpréter.

\subsection{Pas de code dormant}

\begin{invariant}
Les fonctionnalités retirées sont supprimées, non commentées ou laissées non
référencées. L'historique Git est l'archive.
\end{invariant}

Cette règle est inconfortable mais saine : elle interdit de conserver du code
\og au cas où \fg, qui est la principale source d'ambiguïté dans un dépôt. Elle
a conduit à supprimer des milliers de lignes — connecteurs, schémas, composants
d'interface — et à assumer que Git garderait la trace.

\section{Synthèse}

\begin{table}[H]
\centering
\small
\caption{Décisions métier verrouillées lors de la phase de cadrage, et leur
traduction en règle technique.}
\label{tab:decisions-metier}
\begin{tabularx}{\textwidth}{@{}c X X@{}}
\toprule
\textbf{\#} & \textbf{Décision} & \textbf{Conséquence technique} \\
\midrule
1 & L'acheteur est l'assureur français, pilote au poste de souscription &
  Interface orientée dossier unique, pas portefeuille \\
2 & Le livrable principal est un contrat JSON, pas une interface &
  Schéma Pydantic versionné comme artefact central \\
3 & La jointure WFS $\times$ BDNB est le périmètre ; Copernicus et
  Open-Meteo hors v1 &
  Suppression des connecteurs climatiques de la version 1 \\
4 & La transaction est l'adresse unique &
  \api{POST /diagnostic/adresse} comme primitive, le lot en dérivé \\
5 & 13 aléas, certification sur 4 d'entre eux &
  Carte de couches en configuration, pas en logique \\
6 & RGA étiqueté communal &
  Résolution forcée au niveau du schéma \\
7 & Sismique et radon communaux par nature &
  Résolution permanente \code{commune-level} \\
8 & BDNB verbatim &
  Test de conformité clé par clé \\
9 & Score composite mort &
  Schéma qui refuse la sérialisation de champs de score \\
10 & Bandes D03 en présentation seulement &
  Coloration côté frontend, hors contrat \\
11 & Le champ de résolution est le signal de qualité &
  Tests de dégradation aussi rigoureux que les tests nominaux \\
12 & La carte reste, mais montre la jointure &
  Empreinte colorée par le test, panneau de vulnérabilité \\
13 & Le géocodage est un composant produit &
  Tests dédiés : arrondissements, seuil de confiance, saturation \\
\bottomrule
\end{tabularx}
\end{table}

% =====================================================================
\chapter{Analyse de l'existant et recentrage}
\label{ch:existant}
% =====================================================================

Ce chapitre raconte la décision la plus importante du stage, et la moins
confortable : supprimer une grande partie de ce qui avait été construit.

\section{État de départ : un prototype fonctionnel}

Au début du stage, le dépôt contenait déjà une plateforme climatique pour
l'immobilier français. Elle interrogeait quatre sources publiques — Géorisques,
BDNB, Copernicus, Open-Meteo — et présentait les résultats dans un tableau de
bord. Techniquement, cela fonctionnait : les connecteurs répondaient, les
données s'affichaient, une couche de notation produisait un score, une
visualisation 3D extrudait les bâtiments.

La couche géométrique la plus difficile était même déjà écrite : un parseur GML
3.2 fait main et un test point-in-polygon par lancer de rayon, sans dépendance
à \texttt{shapely}.

\section{Le diagnostic : quatre symptômes, une cause}

En auditant le prototype, quatre problèmes sont apparus.

\subsection{La redondance avec les portails sources}

C'est le problème central. Pour chaque aléa, le prototype affichait ce que
Géorisques affiche déjà, à la même résolution, gratuitement. Pour chaque
bâtiment, il affichait ce que la BDNB affiche déjà. Le test appliqué à chaque
fonctionnalité a été formulé ainsi :

\begin{quote}
\itshape
L'acheteur pourrait-il obtenir ceci directement depuis la source ?
\end{quote}

Si la réponse est oui, la fonctionnalité ne justifie pas son existence. Un
agrégateur qui reproduit ses sources n'apporte pas de valeur : il apporte une
dépendance supplémentaire.

\subsection{Le score composite}

Le prototype produisait un score global par bâtiment, agrégé à partir des
aléas. C'est intuitivement satisfaisant — un chiffre, facile à comparer, facile
à afficher. Mais c'est précisément ce que le métier refuse (chapitre~\ref{ch:metier}) :

\begin{itemize}[leftmargin=1.4em]
  \item un score tiers engage la responsabilité de son producteur ;
  \item un score opaque est soit ignoré, soit reconstruit par l'actuaire ;
  \item un score mélange des aléas non comparables (inondation et sismicité ne
        se compensent pas) ;
  \item sous Solvabilité~II, une décision doit être justifiable par des faits,
        pas par un indice.
\end{itemize}

\subsection{Le tableau de bord climatique}

Copernicus et Open-Meteo alimentaient un panneau de variables climatiques. Le
problème n'était pas la donnée, mais sa forme : les valeurs exposées étaient un
\emph{indice d'exposition de 0 à 100}, dérivé, et non des variables
climatiques en unités réelles (jours de canicule par an, degrés Celsius,
mètres par seconde). Un indice 0--100 n'est pas une donnée climatique : c'est
un jugement déguisé. Il tombait sous le coup de la règle sur les scores.

\subsection{La décoration 3D}

Les extrusions 3D des bâtiments n'apportaient aucune information : elles
montraient une forme sans dire si le bâtiment était exposé. C'était un coût de
performance et de complexité pour zéro contenu.

\subsection{La cause commune}

Ces quatre symptômes ont une racine unique : \textbf{le produit n'était pas
organisé autour de la seule chose qu'il faisait et que personne d'autre ne
faisait} — la jointure entre exposition et vulnérabilité, à l'échelle du
bâtiment. L'effort était réparti sur quatre connecteurs, une couche de
notation, un tableau de bord et de la 3D, plutôt que concentré sur cette
jointure.

\section{La décision de recentrage}

Le 25 août 2026, une session de cadrage a produit une réorientation complète,
documentée dans \fichier{docs/workflow/concept.md} et
\fichier{01-brainstorm-notes.md}. Elle tient en cinq points.

\begin{enumerate}[leftmargin=1.6em]
  \item \textbf{Le livrable est un contrat de données.} L'interface devient un
        outil de vente et de démonstration, pas le produit.
  \item \textbf{Le score composite est mort.} Faits et provenance uniquement.
  \item \textbf{Le périmètre de la v1 est la jointure Géorisques $\times$
        BDNB.} Copernicus et Open-Meteo n'ont pas d'usage assureur en v1.
  \item \textbf{Le champ de résolution devient le signal de qualité.} Puisque
        le score disparaît, c'est lui qui porte l'intégrité du contrat.
  \item \textbf{Le code mort est supprimé, pas mis en sommeil.}
\end{enumerate}

\section{L'exécution : supprimer}

La conséquence la plus visible du recentrage est une phase de suppression.

\begin{table}[H]
\centering
\small
\caption{Suppressions de la phase de recentrage, par catégorie.}
\label{tab:suppressions}
\begin{tabularx}{\textwidth}{@{}l X@{}}
\toprule
\textbf{Catégorie} & \textbf{Éléments supprimés} \\
\midrule
Connecteurs backend &
  \fichier{copernicus.py}, \fichier{era5\_land.py}, \fichier{flood.py},
  \fichier{fire\_danger.py}, \fichier{heat\_stress.py}, \fichier{wind\_risk.py},
  \fichier{seasonal.py} \\
\addlinespace
Routes backend &
  \fichier{climate.py}, routes \api{/diagnostic/copernicus/*} \\
\addlinespace
Schémas &
  \fichier{zone\_assessment.py}, schéma \code{RecommandationsIA} et le champ
  \code{recommandations} \\
\addlinespace
Scoring &
  \fichier{risk\_model.py} et toutes ses références, \fichier{app/scoring/} \\
\addlinespace
Dépendances &
  \code{cdsapi}, \code{xarray}, \code{netCDF4} retirés de
  \fichier{requirements.txt} \\
\addlinespace
Tests &
  \fichier{test\_climate\_dashboard.py}, \fichier{test\_copernicus\_*.py} (4),
  \fichier{test\_trajectoire\_brute.py} \\
\addlinespace
Composants frontend &
  \fichier{CopernicusPanel.tsx}, \fichier{CopernicusStatusBanner.tsx},
  répertoire \fichier{jumeau/} (moteur de scène, LIDAR, recommandations),
  \fichier{Artisans.tsx} \\
\addlinespace
Styles &
  \fichier{copernicus.css}, \fichier{jumeau.css}, \fichier{artisans.css} \\
\bottomrule
\end{tabularx}
\end{table}

Deux règles de méthode ont été respectées pendant ces suppressions. D'abord,
elles constituent un \emph{change set dédié} : on ne supprime pas du code au
milieu d'une tâche de fonctionnalité, car cela rend la revue impossible. Ensuite,
le critère n'était pas \og c'est ancien \fg{} mais \og c'est inatteignable \fg :
avant de supprimer, on vérifie qu'aucune route, aucun import, aucune référence
ne pointe vers l'élément. Un code partagé et vivant, comme le module de
recommandations du répertoire \fichier{jumeau/}, a été \emph{déplacé} plutôt que
supprimé.

\section{Le compromis assumé : ce qui n'a pas été supprimé}

Deux éléments ont été conservés malgré la logique de recentrage, et il faut le
dire.

\begin{description}[leftmargin=0em,style=nextline]
  \item[La carte.] Elle n'est pas de la décoration : le souscripteur fait du
  jugement spatial. Mais elle a été ré-orientée : au lieu d'afficher les
  couches sources, elle affiche la jointure. L'empreinte du bâtiment est
  colorée par le résultat du test géométrique, accompagnée d'une fiche de
  vulnérabilité.
  \item[Les projections climatiques.] La direction produit initiale les
  excluait de la v1, mais la non-stationnarité (chapitre~\ref{ch:metier}) est
  l'argument de fond du produit. Les projections sont donc revenues, sous une
  forme différente : non plus comme un tableau de bord de variables, mais comme
  une extension honnête du contrat, avec de vraies unités et une étiquette de
  résolution.
\end{description}

\section{Enseignement}

La leçon de ce chapitre est méthodologique. Un projet technique n'est pas
\og bien avancé \fg{} parce qu'il contient beaucoup de code. Il est bien avancé
quand il devient difficile de lui retirer quelque chose sans le casser. Le
recentrage a réduit le périmètre fonctionnel de manière significative et
augmenté la valeur du livrable — parce que la valeur n'était pas dans le
nombre de fonctionnalités, mais dans l'exactitude de la jointure et dans la
traçabilité de chaque fait.

\begin{figure}[H]
\centering
\begin{tikzpicture}[
  scale=0.95, transform shape,
  box/.style={draw,rounded corners=2pt,align=center,inner sep=5pt,font=\small},
  removed/.style={box,draw=typhoonRed,fill=typhoonRed!6,text width=2.6cm},
  kept/.style={box,draw=typhoonGreen,fill=typhoonGreen!8,text width=2.6cm},
  core/.style={box,draw=typhoonBlue,very thick,fill=typhoonBlue!10,text width=3.4cm}
]
\node[removed] (c1) at (0,2.4) {Connecteurs climatiques};
\node[removed] (c2) at (0,1.2) {Couche de notation};
\node[removed] (c3) at (0,0) {Tableau de bord climat};
\node[removed] (c4) at (0,-1.2) {Extrusions 3D décoratives};
\node[removed] (c5) at (0,-2.4) {Rapport narratif IA};

\node[core] (core) at (5.6,0) {\textbf{La jointure}\\[3pt]
  Aléa Géorisques $\times$ Vulnérabilité BDNB\\[3pt]
  {\footnotesize par bâtiment, avec provenance}};

\node[kept] (k1) at (11.2,1.8) {Carte de la jointure};
\node[kept] (k2) at (11.2,0.6) {13 aléas normalisés};
\node[kept] (k3) at (11.2,-0.6) {Contrat JSON canonique};
\node[kept] (k4) at (11.2,-1.8) {Moteur de scénarios TRI};

\foreach \x in {c1,c2,c3,c4,c5} \draw[-{Latex[length=2mm]},typhoonRed,dashed] (\x) -- (core);
\foreach \x in {k1,k2,k3,k4} \draw[-{Latex[length=2mm]},typhoonGreen] (core) -- (\x);

\node[font=\small\itshape,color=typhoonRed] at (2.6,3.3) {retiré du périmètre};
\node[font=\small\itshape,color=typhoonGreen] at (11.2,3.3) {conservé / construit};
\end{tikzpicture}
\caption{Le recentrage en une image. Tout ce qui ne sert pas la jointure
aléa $\times$ vulnérabilité est retiré du périmètre ou relégué à la couche de
présentation.}
\label{fig:recentrage}
\end{figure}

% =====================================================================
\chapter{Méthodologie d'ingénierie}
\label{ch:methodo}
% =====================================================================

Un stage court impose une discipline : avec une durée limitée, la principale
cause d'échec n'est pas la difficulté technique, c'est la dispersion. Ce
chapitre décrit le cadre méthodologique mis en place, inspiré du développement
piloté par les spécifications (\emph{spec-driven development}), et les raisons
pour lesquelles il a été nécessaire.

\section{Le problème que la méthode résout}

Le projet a une particularité : une part importante du travail a été réalisée
en collaboration avec des agents de programmation assistée par IA. Ce mode de
travail est puissant pour la production de code, mais il introduit un risque
spécifique : l'agent qui reprend le travail ne connaît pas les décisions
antérieures, et peut \emph{défaire} un choix de conception sans s'en rendre
compte.

Deux exemples réels illustrent ce risque :

\begin{itemize}[leftmargin=1.4em]
  \item Un agent a réintroduit un niveau d'eau dépendant de la météo du jour.
        Conséquence : par temps sec, tous les scénarios affichaient 0~m, ce qui
        se lit comme une panne. Le comportement avait pourtant été
        explicitement corrigé.
  \item Un agent a utilisé une boîte englobante à l'échelle du quartier pour
        interroger le WFS. Mesuré à Nantes : la boîte se remplissait de
        polygones situés à 2~km, et le polygone contenant réellement l'adresse
        n'était jamais renvoyé dans la page de résultats. Résultat : \og hors
        TRI \fg{} sur un quai inondable.
\end{itemize}

Dans les deux cas, le code \emph{fonctionnait} au sens où il s'exécutait sans
erreur. Le défaut était sémantique. Une suite de tests ne suffit pas si les
tests eux-mêmes ne sont pas ancrés dans des décisions explicites.

\section{Le triptyque documentaire}

Trois documents structurent le projet, dans un ordre imposé.

\subsection{La constitution}

La \emph{constitution} (\fichier{constitution.md}) énonce les règles non
négociables. Chaque règle porte une étiquette d'application :

\begin{table}[H]
\centering
\small
\caption{Étiquettes d'application des règles constitutionnelles.}
\label{tab:enforcement}
\begin{tabularx}{\textwidth}{@{}l l X@{}}
\toprule
\textbf{Étiquette} & \textbf{Mécanisme} & \textbf{Exemple} \\
\midrule
\texttt{[CI]} & Vérification automatisée &
  Porte de recherche qui échoue si \code{niveau}, \code{score} ou
  \code{recommandations} apparaissent dans un schéma émis \\
\addlinespace
\texttt{[test]} & Assertion dans la suite &
  Le schéma refuse la sérialisation d'un objet portant un score, même si cet
  objet en contient un en interne \\
\addlinespace
\texttt{[review]} & Point de contrôle humain &
  Aucune nouvelle dépendance sans décision enregistrée dans le dépôt \\
\bottomrule
\end{tabularx}
\end{table}

La règle de révision est explicite : \emph{pour changer une règle, on amende la
constitution avant le changement qui la violerait}. Cela force la discussion à
se tenir au niveau du principe, pas au niveau du correctif.

\subsection{La spécification}

La spécification (\fichier{docs/workflow/spec.md}) décrit le produit en
exigences fonctionnelles numérotées, \frref{01} à \frref{28}, chacune assortie
de critères d'acceptation vérifiables. Par exemple :

\begin{description}[leftmargin=0em,style=nextline]
  \item[\frref{03}] Rejeter les géocodages de faible confiance avant tout appel
  aval : score $< 0{,}4 \rightarrow$ HTTP 422 \code{adresse\_ambigue} avec
  \code{label\_propose}. \emph{Critère :} une requête ambiguë renvoie 422 et
  aucun appel BDNB ou Géorisques n'a lieu (vérifié par compteurs de mocks).
  \item[\frref{14}] Ne jamais inférer une exposition depuis une source en
  échec : \code{present=true} uniquement depuis un test réussi,
  \code{false} depuis un test négatif, \code{null} si inconnu.
  \item[\frref{26}] Cache de résultats en mémoire, clé sur le triplet de
  géocodage normalisé (libellé + latitude/longitude arrondies à 5 décimales
  $\approx$ 1~m + code INSEE) et sur la version de schéma.
\end{description}

Les critères d'acceptation sont rédigés de manière à être directement
traduisibles en tests. C'est volontaire : une exigence sans test associé est un
souhait, pas une exigence.

\subsection{Le plan}

Le plan (\fichier{docs/workflow/plan.md}) découpe la spécification en tâches
T001 à T018, chacune avec ses dépendances, son estimation, ses points de
contact constitutionnels, ses critères d'acceptation et — point important —
\emph{les tests à écrire en premier}.

\begin{lstlisting}[caption={Extrait du plan : la tâche T008, porte de
certification de la v1.},label={lst:plan-t008}]
### T008 - Certification known-answer suite (v1 sign-off gate)
- Depends on: T004, T007
- Spec reference: FR-11
- Estimate: 4-8 h (fixture collection dominates)
- Acceptance criteria:
  - Recorded fixtures: known-inside and known-outside address per certified
    hazard (inondation/PPRI, mouvement_terrain, feu_foret, avalanche)
  - Inside -> present=true + resolution="per-building"; outside -> present=false
- Tests to write first:
  - test_certified_inside_<hazard> x4
  - test_certified_outside_<hazard> x4
\end{lstlisting}

\section{Le flux de travail}

Le flux impose un ordre qui n'est pas négociable :

\begin{center}
\begin{tikzpicture}[
  scale=0.76, transform shape,
  node distance=6mm,
  stepbox/.style={draw=typhoonBlue,fill=typhoonBlue!6,rounded corners=2pt,
               align=center,inner sep=5pt,font=\small,text width=2.5cm},
  arrow/.style={-{Latex[length=2mm]},thick,typhoonBlue}
]
\node[stepbox] (s1) {Brainstorm\\{\scriptsize décisions brutes}};
\node[stepbox,right=of s1] (s2) {Concept\\{\scriptsize questions résolues}};
\node[stepbox,right=of s2] (s3) {Spec\\{\scriptsize FR-01..FR-28}};
\node[stepbox,right=of s3] (s4) {Plan\\{\scriptsize T001..T018}};
\node[stepbox,right=of s4] (s5) {Test\\{\scriptsize commit qui échoue}};
\node[stepbox,right=of s5] (s6) {Code\\{\scriptsize commit qui passe}};
\foreach \a/\b in {s1/s2,s2/s3,s3/s4,s4/s5,s5/s6}
  \draw[arrow] (\a) -- (\b);
\end{tikzpicture}
\end{center}

Deux règles de ce flux méritent d'être soulignées.

\begin{enumerate}[leftmargin=1.6em]
  \item \textbf{Un changement de contrat passe d'abord par la constitution,
        puis par la spécification — jamais l'inverse.} Si une fonctionnalité
        nécessite de violer un invariant, ce n'est pas la fonctionnalité qu'on
        ajuste en premier, c'est le principe.
  \item \textbf{Les tests existants ne sont jamais modifiés pendant une tâche
        de fonctionnalité.} Un test jugé erroné est signalé et corrigé dans un
        change set séparé. Sans cette règle, il est toujours possible de faire
        passer un test en le réécrivant — ce qui vide la suite de son sens.
\end{enumerate}

\section{Convention de commit}

Les messages suivent une convention préfixée par le domaine :
\code{feat(etape2):}, \code{fix(flood):}, \code{test(SCN-xxx):},
\code{docs:}, \code{chore:}. Pour toute fonctionnalité, deux commits : un
commit \code{test(TXXX)} qui échoue, puis un commit \code{feat(TXXX)} qui
passe. Cette contrainte rend visible, dans l'historique, que le test a bien
précédé le code.

\section{Décisions documentées}

Les décisions qui poseraient problème à un lecteur futur font l'objet de
documents courts. Trois exemples du projet :

\begin{table}[H]
\centering
\small
\caption{Exemples de décisions documentées et leur motivation.}
\label{tab:adr}
\begin{tabularx}{\textwidth}{@{}l X@{}}
\toprule
\textbf{Décision} & \textbf{Document et motivation} \\
\midrule
Suppression du score composite &
  \fichier{concept.md}, OQ-1 : le score engage la responsabilité et se fait
  reconstruire ou refuser par l'actuaire \\
\addlinespace
RGA communal &
  \fichier{concept.md}, OQ-2 : la source vectorielle n'est pas câblée, annoncer
  plutôt que surinterpréter \\
\addlinespace
Licences commerciales &
  \fichier{concept.md}, OQ-3 : Géorisques et BDNB sont en Licence Ouverte 2.0,
  la réutilisation commerciale est permise, l'obligation d'attribution
  coïncide avec la conception de provenance \\
\addlinespace
Mode VFX &
  \fichier{constitution.md} §2.1, amendement du 14 septembre : autoriser un
  rendu cinématique à condition qu'il soit étiqueté non contractuel et
  strictement confiné à la présentation \\
\bottomrule
\end{tabularx}
\end{table}

\section{Outils}

\begin{table}[H]
\centering
\small
\caption{Outils et pratiques du projet.}
\label{tab:outils}
\begin{tabularx}{\textwidth}{@{}l X@{}}
\toprule
\textbf{Domaine} & \textbf{Choix} \\
\midrule
Gestion de version & Git, une tâche = une branche = une PR, messages préfixés \\
\addlinespace
Backend & Python 3.11+, FastAPI, Pydantic v2, \code{httpx} en asynchrone,
  \code{uvicorn} \\
\addlinespace
Frontend & TypeScript, React, Vite, Mapbox GL JS, Vitest \\
\addlinespace
Tests & \code{pytest} + \code{pytest-asyncio} (backend), Vitest (frontend) \\
\addlinespace
Lint / format & \code{ruff} (backend), \code{eslint} + \code{tsc -b} (frontend) \\
\addlinespace
Déploiement & Vercel (mono-projet : statique + fonction Python serverless) \\
\addlinespace
Documentation & Markdown dans le dépôt, lu par les humains et par les agents \\
\bottomrule
\end{tabularx}
\end{table}

\section{Apport de la méthode}

Le bénéfice concret de cette organisation a été de rendre les décisions
\emph{explicites et opposables}. Lorsqu'un comportement a été corrigé, la
correction est inscrite à la fois dans la constitution (le principe), dans la
spécification (l'exigence), dans le plan (la tâche) et dans la suite de tests
(l'assertion). Un agent — ou un humain — qui reprend le code trouve la règle
avant de trouver le piège.

C'est ce qui a permis de stabiliser un projet dont une part du code a été
écrite en collaboration avec des agents, sans que le produit ne dérive.

% =====================================================================
\chapter{Les sources de données et leur intégration}
\label{ch:donnees}
% =====================================================================

\typhoon{} n'invente aucune donnée : il assemble des sources publiques. La
qualité du produit dépend donc entièrement de la qualité de cette intégration.
Ce chapitre décrit chaque source, ce qu'elle apporte, ses contraintes
juridiques et techniques, et les pièges rencontrés.

\section{Vue d'ensemble}

\begin{table}[H]
\centering
\footnotesize
\caption{Les sources de données intégrées, leur rôle et leur niveau de
résolution.}
\label{tab:sources}
\begin{tabularx}{\textwidth}{@{}l l l X@{}}
\toprule
\textbf{Source} & \textbf{Rôle} & \textbf{Accès} & \textbf{Résolution} \\
\midrule
Géorisques (WFS) & Périmètres d'aléas & OGC WFS 2.0, GML 3.2 &
  \code{per-building} par test point-in-polygon \\
Géorisques (REST) & Statut communal, arrêtés CatNat & JSON &
  \code{commune-level} \\
BDNB (CSTB) & Vulnérabilité du bâtiment & REST JSON &
  \code{per-building} (identifiant bâtiment) \\
Géoplateforme (IGN) & Géocodage & REST JSON &
  Point + code INSEE \\
Classes TRI & Profondeur d'eau réglementaire & REST JSON &
  Par point, par classe de retour \\
Sandre & Tracé hydrographique & REST JSON & Tracé \\
Open-Meteo & Pluie horaire, débit GloFAS & REST JSON & Point, horaire \\
Panoramax & Photographie terrain & REST JSON & Point, orientée \\
Copernicus CDS & Projections climatiques 2050/2100 & API CDS & Grille $\approx$ 10~km \\
Mistral & Rédaction du rapport & REST & — (présentation) \\
\bottomrule
\end{tabularx}
\end{table}

\section{Géorisques : le cœur de la jointure}

\subsection{Ce que la source fournit}

Géorisques expose le référentiel des aléas réglementaires français. Il sert
deux interfaces distinctes, et cette dualité est structurante :

\begin{itemize}[leftmargin=1.4em]
  \item une interface \textbf{REST} par code INSEE, qui donne le statut
        communal : \og la commune est-elle concernée par un risque
        inondation ? \fg{} ;
  \item une interface \textbf{WFS} (Web Feature Service) qui expose les
        géométries vectorielles, donc permet de tester un point précis contre
        un polygone.
\end{itemize}

\subsection{Le test point-in-polygon}

Le WFS renvoie du GML 3.2, un dialecte XML volumineux. Le choix a été d'écrire
le parseur et le test géométrique à la main, sans dépendance à
\texttt{shapely}. Deux raisons : éviter une dépendance de compilation lourde
dans un contexte serverless, et garder la maîtrise complète des cas limites.

Le test utilise un \emph{lancer de rayon} (\emph{ray casting}) : on compte les
intersections d'un rayon partant du point avec les arêtes du polygone. Un
nombre impair signifie que le point est à l'intérieur. Le premier anneau d'un
polygone est l'anneau extérieur ; les suivants sont des trous, qu'il faut donc
soustraire.

\begin{lstlisting}[caption={Signature du test géométrique et gestion des
géométries non polygonales
(\fichier{backend/app/connectors/georisques\_wfs.py}).},label={lst:pip}]
def point_in_polygon(pt: tuple[float, float], polygon: Any) -> bool:
    """Point (lon, lat) dans un Polygon / MultiPolygon GML.

    L'anneau exterieur est le premier ; les suivants sont des trous.
    """

def geometry_intersects_point(
    geom: dict[str, Any],
    pt: tuple[float, float],
    radius_m: float = PROXIMITY_RADIUS_M,
) -> bool:
    """Le point (lon, lat) touche-t-il la geometrie ?

    Polygones -> point-in-polygon strict. Points/lignes -> test de proximite
    (une canalisation a 30 m est un risque pertinent meme si le batiment
    n'est pas dessus).
    """
\end{lstlisting}

\subsection{Les treize aléas}

Le référentiel est normalisé en treize objets aléa, chacun portant les mêmes
champs obligatoires :

\begin{center}
\small
\code{icpe}, \code{inondation}, \code{sismicite}, \code{mouvement\_terrain},
\code{radon}, \code{rga}, \code{cavite}, \code{feu\_foret}, \code{avalanche},
\code{canalisations}, \code{vent\_cyclonique}, \code{ppr}, \code{ssp}
\end{center}

\subsection{La ventilation par type de PPR}

Un piège est apparu ici. L'agrégat \code{ppr} répond à la question \og un plan
de prévention quelconque touche-t-il ce point ? \fg. C'est insuffisant pour un
travail par péril : un bâtiment peut être dans un périmètre PPR sismique sans
être dans un périmètre PPR inondation, et l'agrégat ne fait pas la différence.

Une table de correspondance a donc été ajoutée, qui ventile les huit couches
PPR par type de péril :

\begin{lstlisting}[caption={Ventilation des couches PPR par type de péril
(\fichier{georisques\_wfs.py}).},label={lst:ppr}]
# Ventilation des 8 couches PPR de WFS_LAYER_MAP["ppr"] par type de peril.
# L'agregat "ppr" ci-dessus repond "un PPR quelconque touche-t-il ce point ?"
# -- utile pour une carte generale, mais pas pour un peril precis.
PPR_TYPE_LAYERS: dict[str, list[str]] = {
    "inondation":        ["ms:PPRN_PERIMETRE_INOND", "ms:PPRN_PERIMETRE_SUBMAR"],
    "mouvement_terrain": ["ms:PPRN_PERIMETRE_MVT"],
    "seisme":            ["ms:PPRN_PERIMETRE_SEISME"],
    "avalanche":         ["ms:PPRN_PERIMETRE_AVALANCHE"],
    ...
}
\end{lstlisting}

\subsection{Le piège de la boîte englobante}

C'est le bug le plus instructif du projet. Pour interroger un WFS, il faut
fournir une \emph{boîte englobante} (BBOX). Le choix naturel est de prendre une
boîte \og raisonnable \fg, à l'échelle d'un quartier.

Erreur. Le WFS applique un plafond au nombre de géométries renvoyées
(paramètre \code{count}). Mesuré à Nantes : une boîte large se remplit de
polygones situés à 2~km, et le polygone contenant réellement l'adresse
\emph{n'arrive jamais dans la page de résultats}. Le test conclut alors \og
hors TRI \fg{} sur un quai inondable.

La correction a été de réduire la boîte à l'échelle du point : environ 55~m de
marge, soit \code{POINT\_\allowbreak MARGIN\_\allowbreak DEG}. Tout polygone
contenant le point coupe
nécessairement une boîte centrée sur lui : élargir n'ajoute aucune information
et fait perdre des réponses.

\begin{honesty}
Un résultat faux n'est pas toujours le signe d'un bug visible. Ici, le service
répondait correctement, sans erreur, avec une information fausse. C'est le
scénario le plus dangereux pour un produit de données : la panne silencieuse.
\end{honesty}

\subsection{Le piège du client HTTP partagé}

Autre découverte contre-intuitive : le WFS Géorisques exige un client
\code{httpx.AsyncClient} \emph{dédié}. La passerelle lie la connexion
persistante (\emph{keep-alive}) à un backend particulier. Mélanger les appels
REST et WFS sur le même client provoque des 404 en cascade — sans rapport
apparent avec la cause. Cette contrainte est documentée comme invariant et
vérifiée par un test d'intégration qui enchaîne REST puis WFS.

\subsection{Une réponse vide est ambiguë}

Une réponse WFS vide peut signifier deux choses : \og aucun périmètre ne
contient ce point \fg{} (fait) ou \og le serveur a hoqueté \fg{} (panne). La
distinction est essentielle pour le produit. Le connecteur rejoue donc la
requête une fois avant de conclure à une absence. C'est un compromis explicite :
un surcoût réseau en échange d'une réduction des faux négatifs.

\subsection{Les arrondissements}

Paris, Lyon et Marseille ont un découpage administratif particulier : les
arrondissements ont leur propre code INSEE (75101 à 75120 pour Paris), mais les
endpoints communaux attendent le code de la commune parente (75056). Sans
retraduction, la jointure échoue silencieusement. La correspondance est
implémentée et testée.

\section{La BDNB : la vulnérabilité}

La Base de Données Nationale des Bâtiments, produite par le CSTB, décrit chaque
bâtiment français. Le connecteur récupère environ 139 champs : usage, hauteur,
emprise au sol, matériaux de construction, année de construction, performance
énergétique, et des indicateurs de fiabilité de l'adressage.

Deux points de conception :

\begin{enumerate}[leftmargin=1.6em]
  \item \textbf{Passage verbatim.} Aucun renommage, aucun filtrage. Un test de
        conformité compare l'ensemble des clés reçues et l'ensemble des clés
        émises : ils doivent être identiques.
  \item \textbf{Ambiguïté d'adresse.} Plusieurs bâtiments peuvent partager une
        adresse. La première ligne renvoyée est traitée comme le bâtiment
        principal ; les autres sont conservées verbatim dans
        \code{autres\_batiments\_meme\_adresse[]} — jamais supprimées. Un test
        couvre le cas de trois bâtiments à une même adresse.
\end{enumerate}

\section{Les classes TRI : la profondeur d'eau réglementaire}

Pour la mise en situation sur la carte (chapitre~\ref{ch:scenarios}), il faut
une profondeur d'eau. La source utilisée est le \textbf{Territoire à Risque
d'Inondation} (TRI), issu de la Directive Inondation. Le TRI classe
l'exposition en quatre niveaux correspondant à des périodes de retour :

\begin{table}[H]
\centering
\small
\caption{Les classes TRI officielles et leur signification.}
\label{tab:tri}
\begin{tabular}{@{}l l l@{}}
\toprule
\textbf{Classe} & \textbf{Période de retour} & \textbf{Lecture} \\
\midrule
\code{frequent} & $\approx$ 10 ans & Événement courant \\
\code{moyen} & $\approx$ 100 ans & Événement de référence \\
\code{extreme} & $\approx$ 500 ans & Événement rare \\
\code{faible} & au-delà & Exposition résiduelle \\
\bottomrule
\end{tabular}
\end{table}

Le point crucial : ces classes et leurs profondeurs sont des \textbf{facts
réglementaires}. Elles ne sont jamais inventées ni interpolées. Si le point
n'est pas cartographié, le produit le dit — il ne devine pas une hauteur.

\section{Les sources complémentaires}

\begin{description}[leftmargin=0em,style=nextline]
  \item[Géocodage — Géoplateforme (IGN).] L'adresse saisie est convertie en
  coordonnées, code INSEE, code postal et libellé normalisé, avec un score de
  confiance. Le géocodage est traité comme un composant produit, pas comme une
  utilité : gestion des arrondissements, rejet des scores faibles, gestion de
  la saturation (HTTP 429 avec en-tête \code{retry-after}).
  \item[Hydrographie — Sandre.] Le tracé des cours d'eau, utilisé pour
  l'animation de suivi de l'eau sur la carte.
  \item[Météo — Open-Meteo.] Pluie horaire et débit GloFAS. Utilisée pour la
  \emph{forme} de la montée des eaux, jamais pour son niveau
  (voir chapitre~\ref{ch:scenarios}).
  \item[Photographie — Panoramax.] Une photo de terrain orientée sur le point,
  utile au souscripteur pour un contrôle visuel rapide.
  \item[Copernicus CDS.] Les projections climatiques à 2050 et 2100, en vraies
  unités (jours de canicule par an, degrés, mètres par seconde), pour les deux
  scénarios RCP 4.5 et RCP 8.5.
  \item[Mistral.] La rédaction du rapport narratif. Strictement confinée à la
  couche de présentation : elle ne produit aucun fait et n'entre jamais dans
  le contrat canonique.
\end{description}

\section{Contraintes juridiques et quotas}

\subsection{Licence Ouverte 2.0}

Géorisques et la BDNB sont distribués sous \textbf{Licence Ouverte 2.0}
(Etalab). La réutilisation commerciale est explicitement autorisée, sans
limitation de durée ni de territoire. La seule obligation est
l'\textbf{attribution} : nom de la source et date de dernière mise à jour.

Cette obligation coïncide exactement avec la conception de provenance du
produit : la conformité juridique devient une fonctionnalité, pas une charge.
Chaque bloc de provenance du contrat porte une chaîne d'attribution du type :

\begin{center}
\code{Source Géorisques (BRGM / MTE) — données à jour au 2026-07-01}
\end{center}

Le millésime est stocké en configuration par source et mis à jour lorsque le
fournisseur publie une nouvelle version.

\subsection{Quotas}

Les quotas des fournisseurs sont respectés côté client, car les dépasser
signifie une coupure de service :

\begin{table}[H]
\centering
\small
\caption{Quotas amont et mécanismes de protection implémentés.}
\label{tab:quotas}
\begin{tabularx}{\textwidth}{@{}l l X@{}}
\toprule
\textbf{Fournisseur} & \textbf{Quota} & \textbf{Protection} \\
\midrule
BDNB & 120 req/min/IP, 10\,000 req/mois (offre gratuite) &
  Limitation de débit côté client + compteur budgétaire mensuel \\
\addlinespace
Géoplateforme & 50 req/s/IP &
  Garde de rafale + gestion du 429 avec réattente unique \\
\bottomrule
\end{tabularx}
\end{table}

Un garde budgétaire mensuel protège l'allocation BDNB : lorsque le compteur est
épuisé, les soumissions de lots sont rejetées avec un HTTP 429
\code{budget\_epuise}, tandis que les requêtes unitaires interactives
continuent de fonctionner. Augmenter le débit réel est une action commerciale
(souscrire à l'offre supérieure), pas une décision de code.

\section{Le cache}

Un cache en mémoire à durée de vie limitée (TTL, 24~h par défaut,
configurable) évite de re-interroger les sources pour une même adresse. La clé
est construite sur le triplet de géocodage normalisé et sur la version de
schéma :

\begin{center}
\code{libellé normalisé} + \code{latitude/longitude arrondies à 5 décimales} +
\code{code INSEE} + \code{schema\_version}
\end{center}

L'arrondi à 5 décimales correspond à environ 1~m : deux adresses distinctes ne
partagent donc pas une clé par accident. Un point d'honnêteté : un résultat
servi depuis le cache conserve ses horodatages d'origine. Il indique la date de
la collecte réelle, pas la date de la consultation.

\begin{honesty}
Sur l'instance serverless déployée, ce cache ne persiste pas — le système de
fichiers de la fonction est éphémère. Chaque rapport narratif consomme donc un
appel au modèle, et la prose varie légèrement entre deux requêtes identiques.
En local, le cache rejoue le rapport à l'identique. Cette différence est
documentée dans le README plutôt que masquée.
\end{honesty}

% =====================================================================
\chapter{Architecture et contrat de données}
\label{ch:architecture}
% =====================================================================

\section{Le principe : une seule primitive}

L'architecture repose sur une idée simple : il existe \emph{une} transaction
produit, et toutes les autres surfaces n'en sont que des rendus ou des
répétitions.

\begin{encadre}
\textbf{La transaction unitaire} \api{POST /diagnostic/adresse} prend une
adresse en entrée et renvoie un enregistrement JSON canonique. La carte
l'affiche, le lot la répète, l'export PDF la met en page. Aucune autre surface
n'a le droit de posséder sa propre logique métier.
\end{encadre}

Cette contrainte a une conséquence forte : toute amélioration de la qualité de
la transaction unitaire bénéficie automatiquement à l'ensemble du produit. Elle
interdit aussi la classe de bugs la plus coûteuse pour un contrat de données —
deux implémentations qui divergent lentement.

\section{Vue en couches}

\begin{figure}[H]
\centering
\begin{tikzpicture}[
  scale=0.92, transform shape,
layer/.style={draw,rounded corners=3pt,minimum width=14.6cm,text width=14.6cm,
              minimum height=1.55cm,align=center,font=\small},
  arrow/.style={{Latex[length=2.4mm]}-{Latex[length=2.4mm]},thick,typhoonBlue}
]
\node[layer,fill=typhoonBlue!12,draw=typhoonBlue] (ui) at (0,0)
  {\textbf{Couche de présentation}\\[2pt]
   {\footnotesize React + TypeScript · carte Mapbox GL JS 3D · panneaux
   scénarios · export PDF · mode VFX}};

\node[layer,fill=typhoonGreen!10,draw=typhoonGreen] (svc) at (0,-2.1)
  {\textbf{Couche de service (FastAPI)}\\[2pt]
   {\footnotesize pipeline canonique · lot · cache · quotas · budget ·
   authentification · limitation de débit · rapport narratif}};

\node[layer,fill=typhoonSand,draw=typhoonGrey] (conn) at (0,-4.2)
  {\textbf{Couche connecteurs}\\[2pt]
   {\footnotesize géocodage · georisques REST · georisques WFS ·
   bdnb · flood\_alea (TRI) · hydro · meteo · panoramax · copernicus · mistral}};

\node[layer,fill=white,draw=typhoonGrey!60] (ext) at (0,-6.3)
  {\textbf{Sources externes}\\[2pt]
   {\footnotesize Géoplateforme · Géorisques · BDNB · Sandre · Open-Meteo ·
   Panoramax · Copernicus CDS}};

\draw[arrow] (ui) -- (svc);
\draw[arrow] (svc) -- (conn);
\draw[arrow] (conn) -- (ext);
\end{tikzpicture}
\caption{Architecture en quatre couches. La règle d'or : aucune couche
supérieure n'appelle une source externe directement, et aucun fait ne remonte
sans sa provenance.}
\label{fig:couches}
\end{figure}

\section{Structure du dépôt}

\begin{lstlisting}[caption={Organisation des fichiers du projet (simplifiée).},
label={lst:arbo}]
backend/
  app/api/routes/        diagnostic, flood, hydro, photo, report,
                         geocoding, health
  app/connectors/        bdnb, georisques (REST + WFS), flood_alea (TRI),
                         hydro (Sandre), meteo (Open-Meteo), copernicus,
                         geocoding (BAN), panoramax, mistral
  app/services/          canonical (pipeline unitaire), batch, budget,
                         limits, report_agent, latency_report
  app/schemas/           diagnostic_record (contrat serialise), flood,
                         hydro, photo, report, building_data
  app/core/              config, logging, api_key, rate_limit
  tests/                 24 modules de test hors-ligne

frontend/
  src/routes/            Zone (/zone), ReportPage, Dashboard, Portfolio,
                         WatchlistPage, Login, Home, Contact, Faq
  src/components/        UnifiedMap (Mapbox 3D), ScenarioPanel, RiskConsole,
                         RiskPanel, FloodConsole, LeftPanels, VfxDisclaimer
  src/zone/              moteur : floodAlea (client TRI), floodSim, damageModel,
                         impactModel, hazardSim, windSim, exposure,
                         hydroRoute, hydroLayer, pdf-export, canonicalAdapter
  src/zone/vfx/          vfxInputs (bus d'entree du mode cinematique)
  src/typhoon/           chrome applicatif, authentification, navigation

docs/workflow/           brainstorm, concept, spec, plan, scenario-plan,
                         vfx-spec, vfx-plan
constitution.md          regles non negociables
AGENTS.md                conventions et pieges, lu avant toute modification
\end{lstlisting}

\section{Le contrat canonique}

C'est l'artefact central du produit. Sa structure complète figure en
annexe~\ref{ann:contrat} ; on en présente ici la logique.

\subsection{Une adresse, treize aléas, un bâtiment}

\begin{lstlisting}[style=json,caption={Squelette du contrat canonique
(simplifié).},label={lst:contrat}]
{
  "schema_version": "1.0",
  "adresse": {
    "saisie": "Quai de la Rapee, 75012 Paris",
    "normalisee": "...",
    "citycode": "75056", "postcode": "75012", "city": "Paris",
    "lat": 48.848, "lon": 2.370,
    "geocode_score": 0.95
  },
  "aleas": [                    // 13 objets, referentiel complet
    {
      "code": "inondation",
      "libelle": "Inondation",
      "present": true,          // true | false | null (inconnu)
      "present_commune": true,
      "zonage": "Dans un perimetre PPR inondation",   // texte verbatim
      "hauteur_eau_m": 1.5,
      "per_building": {         // present uniquement si teste geometriquement
        "method": "point-in-polygon",
        "radius_m": null,
        "count": 2
      },
      "source": "georisques",
      "url_detail": "https://www.georisques.gouv.fr/...",
      "erreur": null,
      "resolution": "per-building"
    }
  ],
  "bdnb": {
    "donnees": { /* ~139 champs verbatim */ },
    "_source": {
      "provider": "BDNB", "url": "https://api.bdnb.io",
      "attribution": "Source BDNB (CSTB) - donnees a jour au <millesime>",
      "recuperee_le": "2026-08-25T10:00:00+00:00"
    }
  },
  "georisques_source": { /* provider, url, attribution, recuperee_le */ },
  "erreurs_partielles": [],
  "genere_le": "2026-08-25T10:00:01+00:00",
  "avertissement": "Ce rapport agrege les donnees publiques Georisques..."
}

// INTERDIT dans ce schema : niveau / bandes D03, tout score numerique,
// recommandations, copernicus, tout jugement composite. (FR-18, const. §2)
\end{lstlisting}

\subsection{La sémantique du champ \code{present}}

La valeur \code{present} est à trois états, et cette trichotomie porte une
partie de la valeur du produit :

\begin{table}[H]
\centering
\small
\caption{Les trois états de \code{present} et leur signification.}
\label{tab:present}
\begin{tabularx}{\textwidth}{@{}l l X@{}}
\toprule
\textbf{Valeur} & \textbf{Signification} & \textbf{Origine} \\
\midrule
\code{true} & Le point est concerné &
  Un test géométrique réussi, une proximité mesurée, ou un statut communal
  positif \\
\addlinespace
\code{false} & Le point a été testé et n'est pas concerné &
  Un test géométrique négatif, ou un statut communal négatif \\
\addlinespace
\code{null} & On ne sait pas &
  La source a échoué. Le champ \code{erreur} est rempli et une entrée est
  ajoutée à \code{erreurs\_partielles} \\
\bottomrule
\end{tabularx}
\end{table}

La règle \frref{14} interdit formellement d'inférer une valeur booléenne depuis
une source en échec. Un test de propriété le vérifie : quel que soit le mode de
panne injecté (délai d'attente, erreur 5xx, charge malformée, réponse vide), le
résultat est \code{null} avec une erreur renseignée — jamais \code{true} ou
\code{false}.

\begin{honesty}
C'est la règle qui distingue un produit de données honnête d'un produit
\og qui a l'air de marcher \fg. Un aléa absent parce que la source a échoué
ressemble exactement à un aléa absent parce qu'il n'existe pas. Sans la
trichotomie, l'utilisateur conclut \og pas de risque \fg{} là où il devrait lire
\og je ne sais pas \fg.
\end{honesty}

\subsection{La règle de sérialisation des champs optionnels}

Un détail qui a son importance pour un consommateur automatique : les champs
spécifiques à un aléa sont \emph{toujours présents} dans le JSON, à
\code{null} lorsqu'ils ne s'appliquent pas. Une clé n'est jamais absente. Cela
évite au client d'avoir à distinguer \og clé absente \fg{} de \og valeur
nulle \fg, distinction qui est une source classique de bugs d'intégration.

Une exception documentée : l'objet \code{per\_building} n'existe que si le
point a effectivement été testé géométriquement. Sa présence est elle-même
porteuse d'information.

\subsection{Le schéma comme garde-fou}

L'invariant \og pas de score \fg n'est pas une convention de style : il est
imposé par le schéma Pydantic. Un test vérifie qu'un objet contenant
temporairement un score en interne se sérialise \emph{sans} les clés
interdites. Le schéma refuse activement les champs \code{niveau}, les scores
numériques, \code{recommandations} et \code{copernicus}.

\section{Le pipeline unitaire}

\begin{figure}[H]
\centering
\begin{tikzpicture}[
  scale=0.88, transform shape,
  node distance=6mm,
  st/.style={draw=typhoonBlue,fill=typhoonBlue!6,rounded corners=2pt,
             align=center,inner sep=5pt,font=\footnotesize,text width=2.15cm},
  dec/.style={draw=typhoonAmber,fill=typhoonAmber!8,rounded corners=2pt,
             align=center,inner sep=5pt,font=\footnotesize,text width=2.15cm,
             diamond,aspect=1.9},
  arrow/.style={-{Latex[length=2mm]},thick,typhoonBlue}
]
\node[st] (addr) {Adresse saisie};
\node[st,right=of addr] (geo) {Geocodage\\{\scriptsize score $\geq$ 0,4}};
\node[st,right=of geo] (cache) {Cache ?\\{\scriptsize TTL 24 h}};
\node[st,right=of cache] (fan) {Appels paralleles\\{\scriptsize WFS + BDNB + REST}};
\node[st,right=of fan] (norm) {Normalisation\\{\scriptsize 13 aleas + resolution}};
\node[st,right=of norm] (rec) {Contrat\\{\scriptsize DiagnosticRecord}};

\foreach \a/\b in {addr/geo,geo/cache,cache/fan,fan/norm,norm/rec}
  \draw[arrow] (\a) -- (\b);

\node[dec,below=12mm of geo] (rej) {score $<$ 0,4\\422};
\node[dec,below=12mm of fan] (deg) {source en echec\\present $=$ null};
\draw[arrow,typhoonAmber] (geo) -- (rej);
\draw[arrow,typhoonAmber] (fan) -- (deg);
\draw[arrow,typhoonAmber,dashed] (deg) -- (norm);
\end{tikzpicture}
\caption{Le pipeline unitaire. Les sorties dégradées (en ambre) traversent le
même chemin que les sorties nominales : il n'existe pas de code séparé pour
les cas d'échec.}
\label{fig:pipeline}
\end{figure}

\subsection{Le géocodage d'abord, et comme point de non-retour}

Le géocodage est la première étape et constitue un point de non-retour. Sans
coordonnées résolues, il n'y a rien à joindre : une adresse introuvable ou trop
ambiguë arrête la transaction avec un HTTP 422. C'est le seul cas où le pipeline
échoue plutôt que de se dégrader, et c'est logique — la dégradation suppose
qu'il existe quelque chose à dégrader.

Le seuil de confiance est fixé à 0,4. En dessous, la réponse propose le libellé
le plus proche. Un test vérifie qu'aucun appel BDNB ou Géorisques n'a lieu dans
ce cas (compteurs de mocks).

\subsection{Collecte parallèle}

Les sources sont interrogées en parallèle via une structure de collecte
sécurisée (\code{\_safe\_call}) : chaque appel est isolé, avec son délai
d'attente et son traitement d'erreur. L'échec d'une source ne fait jamais
échouer la transaction.

\subsection{Normalisation}

La normalisation est l'étape où la valeur du produit se construit : elle
transforme les charges brutes hétérogènes en treize objets aléa homogènes,
chacun portant sa provenance et son niveau de résolution. C'est également elle
qui applique les règles forcées : RGA toujours
\code{commune-level-estimate}, sismique et radon toujours
\code{commune-level}.

\section{Les routes de l'API}

\begin{table}[H]
\centering
\footnotesize
\caption{Surface de l'API backend.}
\label{tab:routes}
\begin{tabularx}{\textwidth}{@{}l X@{}}
\toprule
\textbf{Route} & \textbf{Description} \\
\midrule
\api{POST /diagnostic/adresse} & Transaction canonique : adresse $\rightarrow$
  enregistrement complet \\
\api{POST /diagnostic/batch} & Soumission d'un lot ($\leq$ 10\,000 adresses) ;
  renvoie un identifiant \\
\api{GET /diagnostic/batch/\{id\}} & Interrogation de l'état du lot \\
\api{GET /diagnostic/adresse/rapport-pdf} & Relais vers le PDF officiel
  Géorisques (commodité souscripteur) \\
\api{GET /diagnostic/zone/building(s)} & Fiches BDNB autour du point \\
\api{GET /api/flood-alea} & Classes TRI au point (WFS, point-in-polygon) \\
\api{GET /api/hydro}, \api{/api/hydro/extend} & Tracé hydrographique (Sandre) \\
\api{GET /api/meteo} & Pluie horaire + débit GloFAS (Open-Meteo) \\
\api{GET /api/photo} & Photographie Panoramax et crédit \\
\api{GET /api/report}, \api{/api/report/stream} & Rapport narratif (Mistral,
  repli sur gabarit si la clé est absente) \\
\api{GET /api/geocode/search} & Autocomplétion d'adresse (BAN) \\
\api{GET /health}, \api{/health/detailed} & Sonde de vie (non authentifiée) \\
\bottomrule
\end{tabularx}
\end{table}

\section{Sécurité}

\subsection{Authentification par clé}

L'authentification est volontairement simple : une clé statique transmise dans
l'en-tête \code{X-API-Key}, validée contre une liste configurée par variable
d'environnement. Une clé absente produit un 401 ; une clé invalide, un 403. Les
routes de santé sont exemptées.

Un choix explicite : \textbf{les clés ne sont jamais journalisées}. Un test
capture les journaux et vérifie que le matériel de clé est absent. Les adresses,
en revanche, peuvent apparaître dans les journaux applicatifs — elles
constituent l'interface de requête, et les tracer est nécessaire pour le
diagnostic.

\subsection{Défenses en profondeur}

\begin{itemize}[leftmargin=1.4em]
  \item Tous les appels externes portent un délai d'attente explicite. Aucune
        attente non bornée n'existe dans le code.
  \item Une limitation de débit entrante protège l'API \typhoon{} elle-même,
        distincte de la conformité aux quotas amont.
  \item Aucun secret dans le code : toutes les informations d'authentification
        passent par des variables d'environnement, et le fichier
        \fichier{.env.example} documente chaque variable.
\end{itemize}

\section{Le mode dégradé}

Le traitement des pannes est conçu comme une fonctionnalité de premier ordre.
Le tableau~\ref{tab:erreurs} résume le comportement attendu dans chaque
situation.

\begin{table}[H]
\centering
\small
\caption{Matrice des erreurs : ce que le client observe dans chaque situation.}
\label{tab:erreurs}
\begin{tabularx}{\textwidth}{@{}l X@{}}
\toprule
\textbf{Condition} & \textbf{Comportement observable} \\
\midrule
Adresse introuvable & 422 \code{adresse\_non\_trouvee}, avant tout appel source \\
Score de géocodage $<$ 0,4 & 422 \code{adresse\_ambigue} avec libellé proposé \\
Géocodeur indisponible & 429 honoré une fois via \code{retry-after}, puis 502
  \code{geocodage\_indisponible} \\
Une couche WFS en échec & Cet aléa : \code{present=null}, \code{erreur}
  renseignée, entrée dans \code{erreurs\_partielles} — réponse 200 partielle \\
Géorisques entièrement indisponible & Les faits communaux disponibles sont
  conservés ; les faits géométriques passent à \code{null} \\
BDNB indisponible & \code{bdnb=null} + entrée dans \code{erreurs\_partielles} \\
Clé absente ou invalide & 401 / 403, avant tout traitement \\
Échec d'un élément de lot & L'élément est signalé individuellement, le lot
  continue \\
Exception inattendue & Journalisée avec trace, 502 structuré — jamais de 500 HTML \\
\bottomrule
\end{tabularx}
\end{table}

La règle de conception sous-jacente : \emph{le mode dégradé est testé avec la
même rigueur que le chemin nominal}. Une suite de tests dite
\og chaos \fg{} éteint chaque source indépendamment, dans chacun de ses modes
de défaillance, et vérifie que la réponse reste un 200 partiel structuré.

% =====================================================================
\chapter{Le moteur de scénarios interactifs}
\label{ch:scenarios}
% =====================================================================

Un contrat JSON est difficile à vendre. Un souscripteur qui voit un JSON
comprend la donnée, mais il ne \emph{ressent} pas le risque. C'est la raison
d'être de la mise en situation sur la carte : rendre visible, sans mentir, ce
que le contrat décrit.

Ce chapitre raconte la construction de ce moteur — et les trois bugs
sémantiques qui ont dû être corrigés, chacun révélant une confusion
conceptuelle plus profonde qu'un simple défaut de code.

\section{Objectif et référence}

L'objectif est une boucle : choisir une classe de retour $\rightarrow$ l'eau
monte sur le bâtiment sur la carte ; déplacer un curseur horaire $\rightarrow$
le niveau évolue en direct ; appuyer sur \og Play \fg{} $\rightarrow$
l'animation se déroule jusqu'au pic.

La référence visuelle assumée est un projet de démonstration de visualisation
d'inondation (FuseLab), connu pour sa capacité à rendre une montée des eaux
lisible et crédible. L'enjeu était de reproduire cette lisibilité \emph{sans}
enfreindre les invariants du produit : aucune profondeur inventée, aucune
affirmation non étiquetée.

\begin{figure}[H]
\centering
\includegraphics[width=0.92\textwidth]{../docs/screenshots/step2-scenario.png}
\caption{Étape 2 — simulation de scénario. À gauche, les classes TRI
officielles ; à droite, la timeline horaire et le bouton de lecture. L'eau
monte sur le bâtiment par extrusion Mapbox.}
\label{fig:step2}
\end{figure}

\section{Le bus d'état unique}

L'architecture repose sur un bus d'état minimal, partagé par la carte et les
panneaux :

\begin{center}
\code{scenarioKey} (classe TRI) \quad·\quad \code{hazardEvent} (aléa
sélectionné) \quad·\quad \code{riskTimeMin} (heure simulée)
\end{center}

Chaque aléa branche un rendu cartographique distinct — nappe d'eau, cône de
feu, couche de vent — mais tous lisent le même moteur. Aucun panneau ne
recalcule un niveau de façon indépendante : c'est la garantie que la console et
la carte affichent la même profondeur.

\begin{figure}[H]
\centering
\begin{tikzpicture}[
  scale=0.86, transform shape,
  node distance=5mm,
  innode/.style={draw=typhoonGrey!60,fill=typhoonSand,rounded corners=2pt,
             align=center,inner sep=4pt,font=\scriptsize,text width=2.3cm},
  st/.style={draw=typhoonBlue,fill=typhoonBlue!7,rounded corners=2pt,
             align=center,inner sep=4pt,font=\scriptsize,text width=2.3cm},
  eng/.style={draw=typhoonGreen,fill=typhoonGreen!8,rounded corners=2pt,
             align=center,inner sep=4pt,font=\scriptsize,text width=3.1cm},
  outnode/.style={draw=typhoonDeep,fill=typhoonDeep!8,rounded corners=2pt,
             align=center,inner sep=4pt,font=\scriptsize,text width=2.3cm},
  arrow/.style={-{Latex[length=1.8mm]},thick,typhoonBlue}
]
\node[innode] (tri) {Classes TRI\\{\tiny profondeur réglementaire}};
\node[innode,below=of tri] (rain) {Pluie horaire\\{\tiny Open-Meteo}};
\node[innode,below=of rain] (sk) {\code{scenarioKey}\\\code{riskTimeMin}};

\node[eng,right=10mm of rain] (eng) {\textbf{Moteur}\\
  \code{scenarioDepthAt}\\{\tiny le scénario fixe le pic,\\la pluie fixe la forme}};

\node[outnode,right=10mm of eng] (map) {Carte\\{\tiny extrusion d'eau}};
\node[outnode,above=of map] (cons) {Console\\{\tiny mêmes chiffres}};
\node[outnode,below=of map] (pan) {Panneau\\{\tiny vignettes honnêtes}};

\draw[arrow] (tri) -- (eng);
\draw[arrow] (rain) -- (eng);
\draw[arrow] (sk) -- (eng);
\draw[arrow] (eng) -- (map);
\draw[arrow] (eng) -- (cons);
\draw[arrow] (eng) -- (pan);
\end{tikzpicture}
\caption{Un seul bus d'état, un seul moteur. La carte et les panneaux ne
peuvent pas diverger.}
\label{fig:bus}
\end{figure}

\section{Bug 1 : les clés de scénario ne correspondaient pas}

\subsection{Symptôme}

Changer de scénario ne produisait rien : la carte restait sèche.

\subsection{Cause}

Le moteur utilisait des clés de scénario (\code{extreme}, \code{moyen},
\code{frequent}, \code{faible}) tandis que la couche d'interface utilisait une
autre table de correspondance (\code{direct}, \code{west}, \code{east},
\code{offshore}), vestige d'une première version où les scénarios étaient
géographiques. La fonction de résolution recevait donc une clé inconnue et
renvoyait \code{undefined}, silencieusement converti en \code{null} — donc
aucune eau.

\subsection{Correction}

Décision \scn{001} : \textbf{les clés du moteur \emph{sont} les clés
réglementaires TRI}. L'indirection est supprimée, ou réduite à un alias
documenté comme déprécié. Il n'existe plus qu'une seule source de vérité.

\section{Bug 2 : le seuil de dommages masquait la carte}

\subsection{Symptôme}

Même avec une clé correcte, la carte ne peignait rien avant un certain niveau.

\subsection{Cause}

Un seuil unique de 0,3~m servait à la fois à décider \og le bâtiment
subit-il des dommages ? \fg{} et \og faut-il peindre de l'eau sur la carte ? \fg.
Or une nappe de 0,1~m est un fait visuel réel, même si elle n'endommage pas
les équipements. Le seuil de dommages masquait la réalité physique.

\subsection{Correction}

Décision \scn{002} : \textbf{deux seuils, jamais confondus}. Le seuil de
0,3~m gouverne les dommages, les compteurs et les tuiles \og inondé \fg. La
carte, elle, peint dès que la profondeur dépasse une valeur plancher de 2~cm
(\code{MAP\_MIN\_DEPTH\_M}), en dessous de laquelle il ne reste qu'un résidu
numérique.

\begin{honesty}
Le seuil plancher de 2~cm est documenté explicitement : \og à t\textsubscript{0}
l'eau est visible \fg{} signifie \og au-delà de 2~cm \fg. C'est un détail, mais
c'est le genre de détail qui, non documenté, devient un rapport de bug
incompréhensible six mois plus tard.
\end{honesty}

\section{Bug 3 : \og rien ne se passe \fg{} par temps sec}

\subsection{Symptôme}

Sur une adresse réelle à Paris, après correction des deux bugs précédents, le
bouton de lecture restait inerte et changer de scénario ne produisait toujours
rien.

\subsection{Cause — deux causes distinctes}

L'investigation sur données réelles a montré deux problèmes indépendants.

\textbf{Cause A : le niveau dépendait de la météo du jour.} Le niveau d'eau
était calculé à partir de la prévision de pluie. Or ce jour-là à Paris, le
profil de pluie était vide (\code{dry: true}, cumul 0~mm) : la profondeur
valait donc 0 et le bouton était désactivé.

C'est une confusion conceptuelle : \emph{une classe de retour
réglementaire n'est pas une prévision météo}. La classe TRI \og 100 ans \fg{}
existe indépendamment du fait qu'il pleuve aujourd'hui.

\textbf{Cause B : la classe ouverte par défaut n'était pas cartographiée.} Le
scénario sélectionné par défaut était \code{extreme}, mais au point considéré
seule la classe \code{faible} était cartographiée. La résolution renvoyait
\code{null}, donc aucune eau.

\subsection{Corrections}

Décision \scn{004a} — \textbf{le scénario pilote le niveau, la pluie ne pilote
que la forme} :

\begin{center}
\code{scenarioDepthAt(profile, hourIndex, peakM)}
\end{center}

La classe TRI officielle fixe le pic — c'est un fait réglementaire. La
prévision de pluie ne fixe que la \emph{forme} de la courbe. En l'absence de
pluie, une rampe divergée documentée prend le relais et l'interface déclare
explicitement \og montée : hypothèse (aucune pluie prévue) \fg. L'ancienne
fonction dépendante de la météo reste exportée mais marquée dépréciée, avec
l'interdiction explicite de l'utiliser comme source du niveau.

Décision \scn{004b} — \textbf{sélection par défaut = classe la plus intense
réellement cartographiée}. Une fonction bascule sur la classe la plus intense
cartographiée, une seule fois par diagnostic, pour ne pas annuler un choix
explicite de l'utilisateur. Les classes non cartographiées deviennent
\code{disabled} avec une infobulle : un clic sans effet se lit comme un bug.

\begin{invariant}
Hors TRI ou absence de classe au point $\rightarrow$ strictement 0~m. Le pic
n'est \emph{jamais} deviné. La fonction \code{scenarioDepthAt} n'exagère que la
forme de la courbe, jamais son maximum.
\end{invariant}

\section{Le mode VFX et l'isolation des modes}

Un second mode de rendu, dit VFX, a été ajouté ultérieurement
(chapitre~\ref{ch:vfx}). Un point de conception important apparaît ici : le
mode VFX consomme \emph{le même} \code{scenarioDepthAt} que le mode données. Il
est donc impossible que les deux modes affichent deux hauteurs différentes.

De plus, lorsque le mode VFX est actif, la nappe d'eau du mode données est
explicitement désactivée (\code{impact=null}), ce qui évite deux nappes
superposées. Cette isolation est vérifiée par un test dédié.

\section{L'absence n'est pas une panne}

Un principe transversal mérite d'être isolé, car il a été formalisé comme
invariant (\scn{003}) : \textbf{l'absence d'un aléa et la panne d'un service
sont deux faits distincts}.

Le connecteur TRI porte cinq vérités différentes :

\begin{enumerate}[leftmargin=1.6em]
  \item chargement en cours ;
  \item service indisponible (panne) ;
  \item le point est dans un TRI mais aucune classe n'y est définie ;
  \item le point est hors TRI ;
  \item l'appartenance n'a pas pu être vérifiée.
\end{enumerate}

La formulation est \emph{partagée} entre le panneau, la console et le rapport :
une seule fonction produit le texte affiché. Cela évite le scénario classique
où trois composants décrivent le même état de trois façons légèrement
différentes, ce qui donne à l'utilisateur l'impression d'un système incohérent.

\begin{table}[H]
\centering
\small
\caption{Les messages distincts pour l'absence d'exposition.}
\label{tab:absence}
\begin{tabularx}{\textwidth}{@{}l X@{}}
\toprule
\textbf{Situation} & \textbf{Message} \\
\midrule
Hors TRI & \og Point non cartographié · Hors zone TRI \fg{} — un fait \\
\addlinespace
Dans un TRI, sans classe au point & Message distinct, sans chiffre — un fait
  réglementaire \\
\addlinespace
Service indisponible & \og Service indisponible · la cartographie TRI n'a pas
  répondu \fg{} — une panne \\
\bottomrule
\end{tabularx}
\end{table}

\section{Les autres aléas}

\subsection{L'ouragan}

En mode ouragan, l'eau est forcée à \code{null} — il n'y a pas d'inondation —
et une couche météorologique synchronisée à l'heure simulée affiche le vent et
la pluie. La console affiche la rafale et la bande correspondante.

\subsection{Le feu}

Pour l'aléa feu, un cône d'exposition est calculé à partir de la direction
réelle du vent et tracé sur la carte.

\begin{lstlisting}[caption={Construction du cône d'exposition au feu.},
label={lst:fire}]
fireConePolygon(meteo, originLonLat) -> GeoJSON Feature
  - cap     = provenance du vent + 180 degres
  - portee  = bornee par FIRE_REACH_MAX_KM = 1.5 km
  - ouverture heritee de fireConeFrom
\end{lstlisting}

Un écart assumé par rapport au plan initial mérite d'être signalé : la couche
devait être une extrusion 3D, elle a été réalisée en remplissage
(\code{fill}) et contour pointillé. La raison est explicite : \emph{un volume
3D sur un cône d'exposition suggérerait une précision qui n'existe pas}. Le
cône est une zone d'exposition possible, pas une emprise de flammes.

Le cône n'est monté que si l'aléa feu est sélectionné \emph{et} que le
diagnostic indique une exposition au feu. Sans météo, aucun cône n'est tracé et
un message explicite s'affiche.

\section{Synthèse des décisions du moteur}

\begin{table}[H]
\centering
\small
\caption{Décisions verrouillées du moteur de scénarios.}
\label{tab:scn}
\begin{tabularx}{\textwidth}{@{}l X@{}}
\toprule
\textbf{Décision} & \textbf{Contenu} \\
\midrule
\scn{001} & Les clés de scénario du moteur sont les clés réglementaires TRI.
  Pas de table de correspondance locale. \\
\addlinespace
\scn{002} & Deux seuils distincts : dommages (0,3~m) et carte (dès 2~cm). \\
\addlinespace
\scn{003} & Absence $\neq$ panne : cinq vérités distinctes, un texte partagé. \\
\addlinespace
\scn{004a} & Le scénario pilote le niveau, la pluie ne pilote que la forme. \\
\addlinespace
\scn{004b} & Sélection par défaut = classe la plus intense cartographiée,
  une fois par diagnostic. \\
\bottomrule
\end{tabularx}
\end{table}

Ces cinq décisions sont inscrites dans \fichier{AGENTS.md}, le document que
tout contributeur — humain ou agent — doit lire avant de modifier le code.
Chacune est accompagnée de la formulation \og ne jamais réintroduire \fg, avec
la raison historique. C'est la réponse méthodologique au problème posé au
chapitre~\ref{ch:methodo} : un correctif qui n'est pas documenté comme
invariant sera défait.

% =====================================================================
\chapter{Le mode de rendu cinématique (VFX)}
\label{ch:vfx}
% =====================================================================

Ce chapitre décrit une fonctionnalité volontairement contradictoire avec le
reste du projet, et la manière dont la contradiction a été encadrée.

\section{Le problème : la tension}

Le produit est fondé sur l'honnêteté de la donnée : pas de score, pas
d'invention, une provenance sur chaque fait, une étiquette de résolution. Une
démonstration cinématique — caméra qui vole, eau qui monte avec des reflets,
particules, effets — est à l'opposé : elle vise l'impression, pas la preuve.

La tentation était double :

\begin{itemize}[leftmargin=1.4em]
  \item \textbf{commerciale} — une démonstration spectaculaire vend un concept
        mieux qu'un JSON ;
  \item \textbf{technique} — un rendu 3D temps réel est un objet d'ingénierie
        intéressant.
\end{itemize}

Le risque était tout aussi double : introduire un rendu qui \emph{paraît}
factuel alors qu'il est illustratif, et laisser ce rendu contaminer le contrat
de données.

\section{La solution : gouverner la contradiction}

Plutôt que de refuser la fonctionnalité ou de l'accepter telle quelle, le
projet l'a encadrée par un amendement constitutionnel. La constitution a été
portée en version~0.2 le 14 septembre 2026, avec un paragraphe §2.1 dédié.

\subsection{Deux modes obligatoires}

\begin{table}[H]
\centering
\small
\caption{Les deux modes de rendu et leurs contrats respectifs.}
\label{tab:modes}
\begin{tabularx}{\textwidth}{@{}l X X@{}}
\toprule
& \textbf{Mode \code{data}} & \textbf{Mode \code{vfx}} \\
\midrule
Rendu & Enveloppes honnêtes, extrusion Mapbox sur les empreintes réelles &
  Couche personnalisée Three.js sur Mapbox, terrain, particules,
  post-traitement \\
\addlinespace
Défaut & \textbf{oui} & non — activé explicitement \\
\addlinespace
Étiquetage & Aucun (le rendu est factuel) &
  Bandeau permanent \og Simulation visuelle — non contractuelle \fg \\
\addlinespace
Données partielles & Mention d'indisponibilité explicite &
  Étiquette secondaire \og Données partielles — rendu illustratif \fg \\
\addlinespace
Contrat API & Inchangé & \textbf{inchangé} \\
\bottomrule
\end{tabularx}
\end{table}

\begin{invariant}
Aucune route backend, aucun champ de schéma, aucune charge de lot ne peut
transporter de géométrie VFX, de score ou de jugement synthétique. Le mode VFX
est strictement confiné à la présentation.
\end{invariant}

\subsection{Règles d'entrée}

Le mode VFX doit consommer \emph{les mêmes} pilotes que le mode données
lorsqu'ils sont disponibles : l'heure simulée, les profils météorologiques, les
bandes de profondeur TRI, le trajet hydrographique, la hauteur et l'emprise
BDNB. Il peut exagérer l'\emph{apparence} — montée lissée, plan d'eau suivant
le terrain, floraison lumineuse, mouvement de caméra — mais il ne peut pas
revendiquer une précision hydraulique.

Les replis illustratifs sont autorisés \emph{uniquement} en mode VFX et
\emph{uniquement} avec l'étiquette de données partielles.

\subsection{Le bus d'entrée unique}

Un module dédié, \fichier{vfxInputs.ts}, est le seul point d'entrée du rendu
cinématique. La règle est explicite : le rendu VFX passe par ce module et rien
d'autre.

\begin{lstlisting}[caption={Le bus d'entrée du mode VFX et sa règle de repli.},
label={lst:vfxinputs}]
vfxWaterLevel({ rainProfile, hourIndex, triPeak, allowFallback })
  -> { levelM, partial, source: 'forecast' | 'fallback' | 'none' }
\end{lstlisting}

Deux garanties ont été codées dans ce module :

\begin{enumerate}[leftmargin=1.6em]
  \item Le repli illustratif est \textbf{plafonné par la classe TRI} :
        $\min(\text{triPeak}, 1{,}0~\text{m})$. Le rendu ne peut donc jamais
        exagérer au-delà de l'exposition réglementaire réelle.
  \item \textbf{Sans classe TRI, la source vaut \code{none} et le niveau
        vaut 0}, même si le repli est autorisé. On n'invente pas une exposition
        qui n'existe pas.
\end{enumerate}

Ce point est important : le mode cinématique peut embellir la \emph{forme},
jamais créer un \emph{fait}. Une adresse hors zone inondable reste sèche, même
en mode spectaculaire.

\begin{honesty}
Un problème d'incohérence a été détecté et corrigé pendant la conception : la
spécification initiale prévoyait une constante de repli de 1,0~m indépendante
de la classe TRI. Elle était à la fois incohérente — elle pouvait afficher plus
que le risque réel — et inatteignable, puisque le mode données consommait une
autre fonction. Elle a été supprimée : le repli ne porte plus que sur la
\emph{forme} de la courbe.
\end{honesty}

\section{Contraintes d'intégration}

\subsection{Exclusion mutuelle avec la couche météo}

La couche météo Windy impose une projection plate (Mercator sans inclinaison).
La caméra 3D du mode VFX, elle, exige une inclinaison forte. Les deux sont
incompatibles.

\begin{invariant}
La couche Windy et la caméra 3D du mode VFX sont mutuellement exclusives tant
que le mode VFX est actif. Cette exclusivité est vérifiée par un test.
\end{invariant}

Concrètement : à l'activation du mode VFX, la couche Windy est désactivée et
l'inclinaison de la carte est forcée au-delà de 55 degrés. À la sortie, la
politique du mode données reprend la main et l'état précédent est restauré.

\subsection{Le montage et le démontage}

Le mode VFX doit s'activer et se désactiver proprement. La contrainte est
double : aucune fuite de contexte WebGL au changement de route, et aucun
doublon de couche en cas de bascule répétée.

\begin{invariant}
La couche VFX se monte et se démonte proprement lors de la bascule de mode.
Aucune fuite de contexte WebGL au changement de route. Objectif de performance :
au moins 30~images par seconde sur un ordinateur portable de milieu de gamme au
niveau de zoom d'un bâtiment avec le plan d'eau actif, avec dégradation
gracieuse — le post-traitement est désactivé en premier.
\end{invariant}

\section{Étiquetage}

L'étiquetage n'est pas cosmétique : c'est la condition qui rend la
fonctionnalité acceptable.

\begin{itemize}[leftmargin=1.4em]
  \item Le bandeau \og \textbf{Simulation visuelle — non contractuelle} \fg{}
        est affiché en permanence en mode VFX.
  \item Une étiquette secondaire \og \textbf{Données partielles — rendu
        illustratif} \fg{} apparaît lorsque les entrées réelles sont
        incomplètes.
  \item Le mode VFX \textbf{ne doit pas réutiliser le vocabulaire du contrat}
        (\code{resolution}, \code{present}, libellés de bandes) comme s'il
        s'agissait d'un fait sourcé.
\end{itemize}

Le composant d'avertissement ne s'affiche \emph{que} en mode VFX : il n'existe
pas de bandeau dormant en mode données.

\section{État d'avancement}

Toutes les tâches du mode VFX n'ont pas été menées à terme pendant la durée du
stage. Le tableau~\ref{tab:vfx-etat} distingue ce qui est livré de ce qui
reste, sans embellissement.

\begin{table}[H]
\centering
\small
\caption{État d'avancement des tâches du mode VFX.}
\label{tab:vfx-etat}
\begin{tabularx}{\textwidth}{@{}l l X@{}}
\toprule
\textbf{Tâche} & \textbf{État} & \textbf{Contenu} \\
\midrule
\vfx{001} & Livré & Amendement constitutionnel, spécification et plan \\
\vfx{002} & Livré & Bascule de mode, bandeau, mise en page atténuée \\
\vfx{003} & Livré & Bus d'entrée, pilote de niveau d'eau, plafonnement TRI \\
\vfx{006} & Livré & Politique de caméra, exclusion mutuelle Windy \\
\vfx{008} & Livré (partiel) & Isolation des modes, pas de double nappe d'eau \\
\vfx{004} & Restant & Dépendance Three.js et couche personnalisée Mapbox \\
\vfx{005} & Restant & Grille de terrain et nuanceur d'eau \\
\vfx{007} & Restant & Réalisateur caméra (\og Play scenario \fg{}) \\
\vfx{009} & Optionnel & Post-traitement \\
\vfx{010} & Optionnel & Effets secondaires vent et feu \\
\bottomrule
\end{tabularx}
\end{table}

Un point d'honnêteté sur cet état : le niveau d'eau calculé par le mode VFX
n'alimente pas encore de couche de rendu 3D, puisque la tâche \vfx{004} n'est
pas achevée. En revanche, ce niveau n'est pas du code dormant : il alimente
l'étiquette de données partielles du bandeau, ce qui est un usage réel — la
constitution interdit le code non référencé.

\section{Bilan de la démarche}

L'intérêt de ce chapitre n'est pas le rendu cinématique lui-même, mais la
manière dont une tension produit a été traitée. Trois éléments ont rendu la
fonctionnalité acceptable :

\begin{enumerate}[leftmargin=1.6em]
  \item \textbf{un amendement explicite} — la règle a été changée avant le
        code, pas après ;
  \item \textbf{une étanchéité vérifiée par des tests} — le mode VFX ne peut
        pas modifier le contrat ;
  \item \textbf{un étiquetage obligatoire} — le spectateur sait toujours qu'il
        regarde une illustration.
\end{enumerate}

Sans ces trois éléments, la fonctionnalité aurait été soit refusée, soit
introduite en silence — et dans le second cas, elle aurait probablement érodé
la confiance dans l'ensemble du produit.

% =====================================================================
\chapter{Qualité, tests et limites connues}
\label{ch:qualite}
% =====================================================================

\section{La stratégie de test}

La suite de tests repose sur trois principes.

\subsection{Hors-ligne par défaut}

La suite tourne sans accès réseau : les sources sont remplacées par des
\emph{fixtures} capturées et des simulations. Les appels réels à Géorisques ou
à la BDNB constituent un étage de tests séparé.

Cette séparation n'est pas une commodité, c'est une nécessité : une suite qui
dépend de services tiers est lente, instable, et échoue pour des raisons qui
n'ont rien à voir avec le code. Une suite qui échoue sans raison est une suite
que l'on finit par ignorer.

\subsection{La dégradation testée aussi rigoureusement que le nominal}

C'est l'exigence \frref{22} de la spécification, et probablement la plus
structurante. Chaque source est éteinte indépendamment, dans chacun de ses
modes de défaillance : délai d'attente, erreur 5xx, charge malformée, réponse
vide. Dans tous les cas, la réponse attendue est un HTTP 200 partiel structuré,
avec les aléas affectés à \code{present=null} et une erreur renseignée. Jamais
un 500, jamais une exception non gagnée.

\subsection{Les tests d'abord}

Le workflow impose d'écrire le test avant l'implémentation : un commit
\code{test(TXXX)} qui échoue, puis un commit \code{feat(TXXX)} qui passe. Deux
règles complémentaires : les tests existants ne sont jamais modifiés pendant une
tâche de fonctionnalité, et un test jugé erroné est corrigé dans un change set
dédié.

\section{Inventaire de la suite}

\begin{table}[H]
\centering
\footnotesize
\caption{Les modules de test backend et leur objet.}
\label{tab:tests-backend}
\begin{tabularx}{\textwidth}{@{}l X@{}}
\toprule
\textbf{Module} & \textbf{Objet} \\
\midrule
\fichier{test\_post\_canonical\_contract.py} & Forme canonique de bout en bout \\
\fichier{test\_diagnostic\_record\_schema.py} & Champs interdits, optionnels
  toujours présents, énumération de résolution \\
\fichier{test\_georisques\_wfs.py} & Parseur GML, ordre des axes, proximité,
  boîte englobante \\
\fichier{test\_wfs\_integration.py} & Client dédié, absence de contamination \\
\fichier{test\_adresse\_georisques.py} & Jointure adresse $\rightarrow$ aléas \\
\fichier{test\_flood\_alea.py} & Classes TRI, absence $\neq$ panne \\
\fichier{test\_certification.py} & Suite de certification à réponse connue \\
\fichier{test\_geocoding\_hardening.py} & Arrondissements, seuil 0,4, 429 \\
\fichier{test\_api\_key\_auth.py} & 401/403, exemptions, clés hors journaux \\
\fichier{test\_result\_cache.py} & Clé de cache, précision, expiration TTL \\
\fichier{test\_upstream\_throttle.py} & Limitation de débit amont \\
\fichier{test\_budget\_guard.py} & Compteur mensuel, rejet du lot seul \\
\fichier{test\_chaos\_degradation.py} & Matrice de pannes par source \\
\fichier{test\_internal\_batch.py} & Isolation des éléments, équivalence unitaire \\
\fichier{test\_rate\_limit.py} & Limitation entrante \\
\fichier{test\_report\_agent.py} & Rapport narratif et repli gabarit \\
\fichier{test\_hydro.py} & Tracé hydrographique \\
\fichier{test\_panoramax.py} & Photographie et crédit \\
\fichier{test\_health.py} & Sondes de vie \\
\fichier{test\_latency\_report.py} & Percentiles de latence \\
\fichier{test\_api\_diagnostic\_offline.py} & Pipeline hors-ligne \\
\bottomrule
\end{tabularx}
\end{table}

Côté frontend, 21 modules de test couvrent le moteur de scénarios, la boucle
complète, les seuils, les vignettes, le mode VFX et l'isolation des modes.

\begin{table}[H]
\centering
\small
\caption{Volumétrie de la suite de tests au terme du stage.}
\label{tab:tests-vol}
\begin{tabular}{@{}l r l@{}}
\toprule
\textbf{Périmètre} & \textbf{Tests} & \textbf{Exécution} \\
\midrule
Backend (pytest, hors-ligne) & 214 & $<$ 60~s \\
Frontend (Vitest) & 143 & quelques secondes \\
\addlinespace
\emph{Total} & \emph{357} & \\
\bottomrule
\end{tabular}
\end{table}

\section{La suite de certification}

La certification \frref{11} mérite un développement à part. Pour quatre aléas —
inondation/PPRI, mouvement de terrain, feu de forêt, avalanche — des fixtures
ont été capturées pour une adresse \emph{connue à l'intérieur} d'un périmètre
et une adresse \emph{connue à l'extérieur}, chacune vérifiée contre les cartes
officielles du portail Géorisques.

L'attendu est strict :

\begin{itemize}[leftmargin=1.4em]
  \item à l'intérieur : \code{present=true} et
        \code{resolution="per-building"} ;
  \item à l'extérieur : \code{present=false}.
\end{itemize}

Ce n'est pas un test unitaire de plus, c'est la \textbf{porte de signature de
la version~1} : il prouve que le test géométrique dit vrai, pas seulement qu'il
s'exécute. La collecte des fixtures est la partie longue de la tâche.

Un script de capture est conservé dans le dépôt, sous le nom
\fichier{record\_certification\_fixtures.py}, avec les
coordonnées candidates et l'attendu pour chaque cas — par exemple la vallée du
Var à Nice comme cas exposé, et la colline de Cimiez comme cas non exposé.

\begin{honesty}
Cette suite est écrite, mais sa \emph{qualification humaine} — c'est-à-dire la
confrontation visuelle des fixtures aux cartes officielles — était encore en
attente à la fin du stage. Le code est là, les fixtures sont là, la validation
finale par un humain reste à faire. C'est une limite, et il vaut mieux
l'énoncer que la laisser découvrir.
\end{honesty}

\section{Les portes de non-régression}

Outre les tests, deux mécanismes automatiques protègent les invariants.

\subsection{Porte de recherche sur les champs interdits}

Une vérification échoue si les identifiants \code{niveau}, \code{score} ou
\code{recommandations} apparaissent dans un schéma émis. C'est la traduction
automatique de l'invariant \og pas de score \fg.

\subsection{Porte de recherche sur les surfaces supprimées}

Après chaque suppression, une vérification confirme qu'aucune référence ne
subsiste vers les modules retirés — \code{copernicus}, \code{mistral},
\code{lidar}, \code{rapport-narratif}. La règle \og pas de code dormant \fg{}
devient ainsi vérifiable mécaniquement.

\section{Le test d'équivalence du lot}

Un test mérite une mention : il vérifie que la sortie d'un élément de lot est
\emph{identique} à celle d'un appel unitaire pour la même adresse, à
l'horodatage près.

C'est la traduction de l'invariant \og le lot est N fois l'unitaire \fg. Sans ce
test, deux implémentations peuvent diverger lentement : un champ ajouté d'un
côté seulement, un arrondi différent, un ordre d'aléas qui change. Le test rend
cette dérive impossible.

\section{Les limites connues}

Un rapport de stage honnête doit énumérer ce qui ne fonctionne pas ou pas
encore. Le tableau~\ref{tab:limites} recense les limites assumées.

\begin{table}[H]
\centering
\small
\caption{Limites connues du système au terme du stage.}
\label{tab:limites}
\begin{tabularx}{\textwidth}{@{}l X@{}}
\toprule
\textbf{Limite} & \textbf{Description et statut} \\
\midrule
Cache non persistant en production &
  En serverless, le système de fichiers de la fonction est éphémère : le cache
  de rapports ne survit pas. Chaque rapport consomme un appel au modèle.
  Documenté. \\
\addlinespace
Store de lot en mémoire &
  Le magasin de lots est en mémoire, mono-processus, sans persistance. Perdu
  au redémarrage. Suffisant pour valider le schéma, pas dimensionné pour
  \og des centaines de milliers d'adresses \fg. \\
\addlinespace
Authentification simulée &
  Sans configuration Supabase, l'authentification reste en mode démonstration :
  la connexion fonctionne sans comptes réels. Non utilisable en production
  externe. \\
\addlinespace
RGA non vectoriel &
  Le retrait-gonflement des argiles reste communal. Un câblage sur une source
  vectorielle (par exemple les aléas argile du BRGM) demande un amendement
  constitutionnel. \\
\addlinespace
Certification non qualifiée &
  La suite de certification à réponse connue attend sa validation humaine
  contre les cartes officielles. \\
\addlinespace
Mode VFX incomplet &
  Les tâches \vfx{004}, \vfx{005} et \vfx{007} (couche Three.js, terrain,
  réalisateur caméra) ne sont pas livrées. \\
\addlinespace
Calibration de latence &
  Le budget de performance (p95 $\leq$ 10~s) est proposé, non encore calibré
  sur une campagne de mesures réelles. Un harness de mesure existe
  (\fichier{latency\_report.py}). \\
\addlinespace
Smoke test manuel non exécuté &
  La vérification manuelle de bout en bout du moteur de scénarios sur une
  adresse certifiée TRI avec pluie prévue n'a pas été exécutée : elle demande
  un jeton Mapbox et des conditions météo particulières. \\
\bottomrule
\end{tabularx}
\end{table}

\begin{honesty}
Le dernier point de ce tableau est le plus important à assumer. Une
fonctionnalité peut être couverte par des tests automatisés et n'avoir jamais
été vue fonctionner de bout en bout sur une donnée réelle. Ce n'est pas la même
chose. La vérification manuelle reste due.
\end{honesty}

\section{Performance}

Le budget de performance proposé par la constitution §6 est le suivant :

\begin{itemize}[leftmargin=1.4em]
  \item p95 $\leq$ 10~s pour \api{POST /diagnostic/adresse} avec toutes les
        sources saines ;
  \item le mode dégradé ne dépasse pas deux fois ce budget — un échec déclenche
        un repli, pas une tempête de réessais ;
  \item toute régression supérieure à 25~\% bloque la fusion.
\end{itemize}

Ces chiffres sont présentés comme \emph{proposés}, à calibrer par la mesure.
Un point d'attention est identifié : la résolution géométrique interroge
actuellement les couches WFS de façon séquentielle, chacune avec son délai
d'attente et sa réattente. Si les mesures dépassent le budget, la parallélisation
des requêtes de couches est prévue comme tâche conditionnelle (T017) — plutôt
que de relever silencieusement le budget.

\begin{honesty}
Relever un budget de performance pour qu'il corresponde à la mesure est une
manière élégante de ne jamais échouer. Le projet a choisi l'inverse : le budget
est un engagement, et s'il n'est pas tenu, c'est le code qui change.
\end{honesty}

% =====================================================================
\chapter{Déploiement et mise en production}
\label{ch:deploiement}
% =====================================================================

Un service non déployé n'est pas un service. Ce chapitre décrit l'architecture
de déploiement retenue, les obstacles rencontrés, et la vérification qui a été
faite sur l'instance en ligne.

\section{L'objectif : une URL, zéro configuration}

Le choix d'architecture de déploiement a été guidé par une contrainte
d'usage : un utilisateur qui reçoit un lien doit pouvoir l'ouvrir et
l'utiliser. Il ne doit pas avoir à configurer une URL de backend, ni à
comprendre un problème de CORS.

Cela a conduit à un déploiement \textbf{mono-projet} : le frontend statique et
le backend FastAPI servis depuis le même domaine.

\begin{figure}[H]
\centering
\begin{tikzpicture}[
  scale=0.9, transform shape,
  node distance=7mm,
  dom/.style={draw=typhoonDeep,fill=typhoonDeep!8,rounded corners=3pt,
              align=center,inner sep=6pt,font=\small,text width=4.2cm},
  fn/.style={draw=typhoonGreen,fill=typhoonGreen!8,rounded corners=3pt,
             align=center,inner sep=5pt,font=\footnotesize,text width=3.1cm},
  ext/.style={draw=typhoonGrey!70,fill=typhoonSand,rounded corners=3pt,
              align=center,inner sep=5pt,font=\footnotesize,text width=3.1cm},
  arrow/.style={-{Latex[length=2mm]},thick,typhoonBlue}
]
\node[dom] (dom) at (0,0) {\textbf{Domaine unique}\\[2pt]
  {\footnotesize \texttt{typhoon-rose.vercel.app}}};

\node[fn] (static) at (-4.6,-2.4) {Frontend statique\\{\tiny build Vite servi depuis /dist}};
\node[fn] (api) at (0,-2.4) {Fonction Python\\{\tiny \texttt{api/index.py} monte FastAPI}};
\node[fn] (proxy) at (4.6,-2.4) {Proxy Vigicrues\\{\tiny relais same-origin}};

\node[ext] (src) at (0,-4.8) {Sources externes\\{\tiny Géorisques · BDNB · IGN · Open-Meteo}};

\draw[arrow] (dom) -- (static);
\draw[arrow] (dom) -- (api);
\draw[arrow] (dom) -- (proxy);
\draw[arrow] (api) -- (src);
\draw[arrow] (proxy) -- (src);
\end{tikzpicture}
\caption{Architecture de déploiement mono-domaine. Tout passe par une seule
URL : pas de CORS, pas de variable d'environnement côté client.}
\label{fig:deploy}
\end{figure}

\section{Le principe du mono-domaine}

Le fichier de configuration à la racine du dépôt orchestre l'ensemble :

\begin{itemize}[leftmargin=1.4em]
  \item \textbf{build} — installation et compilation du frontend, servi depuis
        le répertoire de sortie ;
  \item \textbf{rewrites} — les chemins \api{/api/*}, \api{/diagnostic/*} et
        \api{/health} sont routés vers la fonction Python qui monte
        l'application FastAPI ;
  \item \textbf{proxy Vigicrues} — un relais vers le service de vigilance
        crues, nécessaire car ce service n'envoie pas d'en-têtes CORS. En
        développement, c'est le proxy de Vite qui joue ce rôle ; en production,
        c'est une règle de réécriture ;
  \item \textbf{fallback SPA} — toute route cliente inconnue est servie par
        l'application, pour que \api{/zone} ou \api{/dashboard} ne renvoient
        pas une erreur 404.
\end{itemize}

Côté client, une règle importante : le frontend appelle le backend
\textbf{en même origine} (URLs relatives) dès que la page n'est pas servie
depuis un hôte local. Aucune URL n'est à configurer, aucun CORS n'est à gérer.

\begin{invariant}
Un \code{VITE\_API\_BASE} qui pointe vers \code{http://127.0.0.1:8000} est
\textbf{ignoré en production} : il désignerait la machine du visiteur, ce qui
produirait un message trompeur du type \og backend inaccessible ? \fg. La
résolution d'URL ignore un override local sur un hôte déployé.
\end{invariant}

\section{Les obstacles rencontrés}

Le déploiement a été la partie la plus itérative du stage. L'historique des
commits en garde la trace, et les enseignements valent d'être consignés.

\begin{table}[H]
\centering
\small
\caption{Obstacles de déploiement et résolutions.}
\label{tab:obstacles}
\begin{tabularx}{\textwidth}{@{}X X@{}}
\toprule
\textbf{Symptôme} & \textbf{Cause et correction} \\
\midrule
\og Could not find a top-level \texttt{app} \fg &
  La fonction Python doit exposer le nom \code{app} \textbf{au niveau racine
  du module}, hors bloc conditionnel. Un diagnostic de démarrage peut rester,
  mais l'affectation doit remonter sans indentation. \\
\addlinespace
\og pattern \texttt{api/index.py} doesn't match any Serverless Functions \fg &
  Message trompeur : en configuration moderne, la fonction est écartée
  silencieusement si le nom racine est absent. Le message ne mentionne pas la
  vraie cause. \\
\addlinespace
\code{FUNCTION\_INVOCATION\_FAILED} (500) &
  Variables d'environnement présentes mais vides. La configuration a été
  durcie pour ignorer les variables vides plutôt que de les traiter comme
  définies. \\
\addlinespace
Le répertoire statique monté au mauvais endroit &
  Une configuration héritée montait le résultat du build sous un sous-chemin.
  Correction par une configuration moderne. \\
\addlinespace
\og Vigilance crues indisponible \fg &
  Sans règle de proxy, \api{/vigicrues/...} tombait dans le fallback SPA et
  renvoyait du HTML au lieu de données. \\
\addlinespace
Routes clientes en 404 &
  Ajout du fallback SPA. \\
\bottomrule
\end{tabularx}
\end{table}

Une leçon se dégage de cet ensemble : \textbf{les messages d'erreur d'une
plateforme de déploiement désignent souvent le symptôme, pas la cause}. Le
message \og le motif ne correspond à aucune fonction \fg{} n'indique pas que le
motif est faux, mais que la fonction a été écartée pour une autre raison. Il a
fallu remonter à la cause racine — la visibilité du nom racine — pour corriger.

\section{Comportements propres au serverless}

Le passage à une fonction serverless introduit des différences qu'il faut
connaître avant une démonstration :

\begin{itemize}[leftmargin=1.4em]
  \item \textbf{Démarrage à froid} — le réveil d'une fonction froide prend de
        5 à 10~s, et un diagnostic complet peut atteindre environ 15~s à froid.
        C'est normal sur l'offre gratuite ; l'astuce opérationnelle consiste à
        appeler la sonde de vie une fois avant une démonstration.
  \item \textbf{Cache non persistant} — le système de fichiers de la fonction
        est éphémère. Chaque rapport narratif consomme donc un appel au modèle,
        et la prose varie légèrement entre deux requêtes identiques.
  \item \textbf{Clé Windy liée au domaine} — l'API météo n'autorise une clé que
        sur les domaines déclarés. Le domaine de production doit y figurer,
        sinon les couches météo échouent. La console nomme alors explicitement
        le domaine à ajouter — un bon exemple de message d'erreur actionnable.
\end{itemize}

\section{Vérification de l'instance en ligne}

Une instance de référence a été déployée et vérifiée le 15 septembre 2026.
Le tableau~\ref{tab:verif} reproduit les contrôles effectués et leurs résultats.

\begin{table}[H]
\centering
\small
\caption{Contrôles de l'instance déployée, vérifiés le 15 septembre 2026.}
\label{tab:verif}
\begin{tabularx}{\textwidth}{@{}l X@{}}
\toprule
\textbf{Contrôle} & \textbf{Résultat} \\
\midrule
\api{/} et \api{/zone} & 200 — application servie (fallback SPA) \\
\api{/health}, \api{/health/detailed} & 200, \code{\{"status":"ok"\}} \\
\api{/api/flood-alea} (Quai de la Rapée) & 200 — \code{in\_tri: true}, classes
  TRI réelles \\
\api{/api/meteo}, \api{/api/hydro}, \api{/api/photo} & 200 \\
\api{/api/geocode/search} & 200 — géocodage réel \\
\api{POST /diagnostic/adresse} (Quai de la Rapée) & 200 en environ 9 à 16~s \\
\api{POST /api/report} & 200, repli non utilisé, prose française \\
\api{POST /api/report/stream} & 200 en environ 1,6~s, avec les événements
  \code{header}, \code{summary}, \code{mitigations}, \code{confidence},
  \code{appendix}, \code{patch}, \code{done} \\
Jeton Mapbox injecté dans le bundle & oui \\
\bottomrule
\end{tabularx}
\end{table}

\section{Adresses de référence}

Pour que la démonstration soit reproductible, trois adresses de référence ont
été identifiées et vérifiées.

\begin{table}[H]
\centering
\small
\caption{Adresses de référence pour la démonstration.}
\label{tab:adresses}
\begin{tabularx}{\textwidth}{@{}>{\raggedright\arraybackslash}p{4.2cm}>{\raggedright\arraybackslash}p{4.4cm}>{\raggedright\arraybackslash}X@{}}
\toprule
\textbf{Adresse} & \textbf{Rôle} & \textbf{Résultat vérifié} \\
\midrule
Quai de la Rapée, 75012 Paris & Démonstration principale, eau sur le bâtiment &
  \code{available: true}, Moyen 0--1~m, Faible 2--3~m, bassin Seine \\
\addlinespace
10 Quai de la Charente, 75019 Paris & Absence honnête : dans un TRI sans classe
  au point &
  \code{in\_tri: true}, 6/13 aléas, fiche BDNB 699~m\textsuperscript{2} \\
\addlinespace
Adresse intérieure hors TRI (ex. Brou, 28300) & Contraste \og hors TRI $\neq$
  jamais inondé \fg &
  \code{available: false}, \code{in\_tri: false} \\
\bottomrule
\end{tabularx}
\end{table}

Le parcours de démonstration suggéré exploite délibérément le contraste : la
première adresse montre le produit en fonctionnement nominal, la deuxième
montre la rigueur — un message distinct, aucun chiffre inventé — et la
troisième montre l'honnêteté sur l'absence.

\section{Alternative de déploiement}

Une seconde option de déploiement est documentée pour le cas où les fonctions
serverless seraient trop lentes : backend sur une plateforme d'hébergement
classique (blueprint prêt dans le dépôt) et frontend statique. Le compromis est
inverse : démarrage plus lent après inactivité (environ 30~s de réveil sur
l'offre gratuite), mais performances d'exécution plus régulières. Cette option
existe parce qu'un choix d'infrastructure est un arbitrage, pas une évidence.

% =====================================================================
\chapter{Résultats et validation}
\label{ch:resultats}
% =====================================================================

\section{Le livrable}

Au terme du stage, \typhoon{} est un service déployé et vérifié en ligne,
composé de :

\begin{itemize}[leftmargin=1.4em]
  \item un backend FastAPI de 70 fichiers Python, 8 routes fonctionnelles et
        10 connecteurs ;
  \item un frontend React/TypeScript de 68 modules, avec une carte Mapbox 3D,
        un moteur de scénarios et un mode de rendu cinématique ;
  \item \textbf{357 tests automatisés} hors-ligne (214 backend, 143 frontend) ;
  \item un contrat de données versionné, dont la conformité est imposée par le
        schéma et vérifiée par des tests ;
  \item une constitution d'ingénierie de 8 sections et un amendement ;
  \item une documentation de travail complète : brainstorm, concept,
        spécification (28 exigences), plan (18 tâches), plan de scénarios
        (10 tâches), spécification et plan VFX.
\end{itemize}

\begin{figure}[H]
\centering
\includegraphics[width=0.95\textwidth]{../docs/screenshots/step1-diagnostic.png}
\caption{Étape 1 — diagnostic. Les 13 aléas Géorisques avec leur badge de
résolution, la fiche BDNB du bâtiment, et une photographie terrain Panoramax
orientée sur le point, sur une carte 3D.}
\label{fig:step1}
\end{figure}

\section{Le contrat en pratique}

Le contrat produit pour l'adresse de démonstration principale illustre la
jointure. On y lit, pour chaque aléa, trois informations indissociables : le
fait, sa source, et sa résolution.

\begin{table}[H]
\centering
\small
\caption{Lecture du contrat pour l'adresse de démonstration principale.}
\label{tab:lecture-contrat}
\begin{tabularx}{\textwidth}{@{}l l l X@{}}
\toprule
\textbf{Aléa} & \textbf{\code{present}} & \textbf{Résolution} & \textbf{Origine du fait} \\
\midrule
Inondation & \code{true} & \code{per-building} &
  Test point-in-polygon réussi contre le périmètre PPR inondation \\
\addlinespace
Sismicité & selon zone & \code{commune-level} &
  Zonage sismique décrétal — résolution maximale existante \\
\addlinespace
Radon & selon commune & \code{commune-level} &
  Classement de potentiel communal \\
\addlinespace
RGA & selon commune & \code{commune-level-estimate} &
  Estimation communale assumée, source vectorielle non câblée \\
\addlinespace
Feu de forêt & \code{true}/\code{false} & \code{per-building} &
  Test géométrique contre le périmètre \\
\bottomrule
\end{tabularx}
\end{table}

Le point important n'est pas la valeur de chaque ligne, mais le fait que la
troisième colonne existe. Un lecteur du contrat sait \emph{toujours} à quel
niveau l'information a été établie, et peut donc décider du crédit à lui
accorder.

\section{Ce que le produit démontre}

\subsection{La jointure existe et elle est vérifiée}

Le résultat central est que la jointure entre exposition et vulnérabilité, à
l'échelle du bâtiment, est établie par un test géométrique réel et non par une
approximation communale. Pour les aléas couverts par un périmètre vectoriel, le
fait porte effectivement sur le bâtiment.

\subsection{Le refus du score est tenable}

La suppression du score composite a été testée jusqu'au bout : le schéma refuse
la sérialisation de champs de score, une porte de recherche échoue si un tel
champ apparaît, et la valeur produit se reconstruit sur d'autres fondations —
la provenance et la résolution. Le produit est démontrable sans son chiffre.

\subsection{Les modes de défaillance sont explicites}

La trichotomie \code{true}/\code{false}/\code{null} et la matrice d'erreurs
font que le produit ne confond jamais \og pas de risque \fg{} et \og je ne sais
pas \fg. C'est peut-être la propriété la plus difficile à démontrer sur un
produit de données, et elle est ici vérifiée par une suite de tests de
dégradation.

\subsection{Le moteur de scénarios relie la donnée au jugement}

La carte ne montre pas une donnée : elle montre la jointure. Et le moteur de
scénarios transforme une classe de retour abstraite (\og 100 ans \fg) en une
montée des eaux lisible, sans jamais inventer un niveau.

\section{Chiffres clés}

\begin{table}[H]
\centering
\small
\caption{Synthèse quantitative du projet.}
\label{tab:chiffres}
\begin{tabularx}{\textwidth}{@{}l l X@{}}
\toprule
\textbf{Indicateur} & \textbf{Valeur} & \textbf{Commentaire} \\
\midrule
Commits & 76 & du 14 août au 15 septembre 2026 \\
\addlinespace
Fichiers Python (backend) & 70 & dont 10 connecteurs, 8 routes, 7 services \\
\addlinespace
Modules TypeScript (frontend) & 68 & dont 20 dans le moteur de scénarios \\
\addlinespace
Modules de test & 45 & 24 backend, 21 frontend \\
\addlinespace
Tests automatisés & 357 & 214 backend, 143 frontend, hors-ligne \\
\addlinespace
Exigences fonctionnelles & 28 & \frref{01} à \frref{28}, chacune avec
  critères d'acceptation \\
\addlinespace
Tâches planifiées & 18 + 10 & plan principal + plan de scénarios \\
\addlinespace
Aléas normalisés & 13 & référentiel Géorisques complet \\
\addlinespace
Aléas certifiés & 4 & inondation, mouvement de terrain, feu, avalanche \\
\addlinespace
Sources externes intégrées & 9 & dont 7 publiques \\
\addlinespace
Sections constitutionnelles & 8 & dont un amendement (§2.1, mode VFX) \\
\addlinespace
Décisions documentées & 13 + 5 & cadrage + moteur de scénarios \\
\bottomrule
\end{tabularx}
\end{table}

\section{Validation}

La validation du travail repose sur quatre niveaux complémentaires.

\begin{description}[leftmargin=0em,style=nextline]
  \item[Tests automatisés.] 357 tests hors-ligne, exécutés en moins d'une
  minute pour le backend. Ils couvrent le chemin nominal, la dégradation, les
  invariants de schéma et l'équivalence unitaire/lot.
  \item[Vérification en ligne.] Une instance déployée, contrôlée endpoint par
  endpoint le 15 septembre 2026 (tableau~\ref{tab:verif}).
  \item[Confrontation au métier.] Les décisions structurantes — refus du score,
  variables brutes plutôt qu'un indice, livraison à la fois interactive et en
  lot — proviennent d'échanges réels avec un actuaire et un professionnel du
  secteur, pas d'une hypothèse de conception.
  \item[Reproductibilité de la démonstration.] Trois adresses de référence
  documentées, avec le résultat attendu pour chacune, permettent de rejouer la
  démonstration à l'identique.
\end{description}

\section{Apports personnels}

Au-delà du livrable, ce stage a produit plusieurs apprentissages que je
considère transférables.

\begin{enumerate}[leftmargin=1.6em]
  \item \textbf{La valeur d'un produit de données tient à sa traçabilité, pas à
        son volume.} Une jointure modeste mais vérifiable vaut mieux qu'une
        agrégation large et approximative.
  \item \textbf{Le refus est une décision d'ingénierie.} Supprimer le score,
        supprimer quatre connecteurs, supprimer des milliers de lignes a
        \emph{augmenté} la valeur du livrable.
  \item \textbf{Une panne silencieuse est plus dangereuse qu'une erreur
        visible.} Le bug de la boîte englobante répondait correctement avec une
        information fausse. C'est ce qui m'a conduit à traiter la trichotomie
        \code{true}/\code{false}/\code{null} comme une exigence de premier rang.
  \item \textbf{Un correctif non documenté sera défait.} Travailler avec des
        agents de programmation rend cette règle vitale : l'invariant doit
        exister dans un document lu avant le code.
  \item \textbf{Les sources publiques ont des pièges non documentés.} Le client
        HTTP partagé, la boîte englobante, l'ordre des axes GML : rien de tout
        cela n'est dans la documentation officielle. Le savoir s'accumule par
        la mesure et doit être écrit.
\end{enumerate}

\section{Impact potentiel}

Il serait malhonnête de présenter \typhoon{} comme un produit prêt à être
vendu. C'est un prototype fonctionnel qui démontre une thèse : la jointure
aléa $\times$ vulnérabilité par bâtiment est réalisable à partir de sources
publiques, avec une traçabilité compatible avec les exigences d'un audit.

Le chemin vers un usage réel passe par des étapes identifiées : qualification
humaine de la certification, calibration de latence, authentification réelle,
persistance du store de lot, et surtout validation de l'intérêt par un pilote
chez un assureur. La valeur du prototype est d'avoir rendu cette discussion
concrète — on ne débat plus de l'idée, on débat d'une URL qu'on peut ouvrir.

% =====================================================================
\chapter{Bilan critique}
\label{ch:bilan}
% =====================================================================

Un rapport de stage qui ne présente que des réussites n'est pas un rapport :
c'est une plaquette. Ce chapitre expose ce qui a été difficile, ce qui a mal
tourné, et ce que je ferais différemment.

\section{Ce qui a bien fonctionné}

\subsection{La remise en question du cadrage}

Le moment le plus productif du stage a été la session de recentrage du
25 août. Elle a consisté à admettre qu'un travail déjà réalisé — quatre
connecteurs, une couche de notation, un tableau de bord — n'apportait pas la
valeur attendue, et à le supprimer.

Cette décision a été facilitée par un critère unique et brutal :
\emph{l'acheteur pourrait-il obtenir ceci directement depuis la source ?} Un
critère unique est plus facile à appliquer qu'une liste de priorités, parce
qu'il ne laisse pas de place à la négociation.

\subsection{La confrontation au terrain métier}

Les échanges avec un actuaire et un professionnel de l'assurance ont
directement réécrit la direction produit. Le refus du score composite, la
préférence pour les variables brutes et l'exigence de provenance viennent de
ces conversations, pas d'une intuition.

La leçon est générale : \textbf{il est moins coûteux de se tromper auprès d'un
utilisateur que de se tromper tout seul pendant deux mois}. Ces échanges ont
coûté quelques heures et ont évité de construire un produit que personne
n'aurait voulu.

\subsection{La discipline documentaire}

Le triptyque constitution / spécification / plan a permis de stabiliser un
projet dont une part importante du code a été écrite en collaboration avec des
agents de programmation. Sans ce cadre, le projet aurait très probablement
dérivé : chaque session d'agent aurait pu défaire silencieusement une décision
antérieure.

\section{Ce qui a mal fonctionné}

\subsection{Les bugs sémantiques répétés}

Trois fois de suite, un comportement \emph{fonctionnel} a dû être corrigé parce
qu'il était \emph{faux} :

\begin{enumerate}[leftmargin=1.6em]
  \item les clés de scénario ne correspondaient pas ;
  \item un seuil unique confondait dommages et affichage ;
  \item le niveau d'eau dépendait de la météo du jour.
\end{enumerate}

Dans les trois cas, le code s'exécutait sans erreur. Le fait que ces bugs
soient apparus \emph{après} la mise en service montre une faiblesse de méthode :
les tests vérifiaient que le code ne plantait pas, pas qu'il disait vrai. Une
partie de la réponse a été la suite de tests de scénarios
(\fichier{scenarioLoop.test.ts}), qui verrouille la chaîne complète de la clé
jusqu'à l'eau dessinée. Mais cette suite est arrivée \emph{après} les bugs.

Ce que je retiens : \textbf{pour un produit de données, un test qui vérifie
l'absence d'exception est presque inutile}. Ce qui compte est un test qui
vérifie une valeur attendue sur une entrée connue.

\subsection{La dérive initiale}

La phase 1 a produit du code qui a dû être supprimé. Avec le recul, le
problème n'était pas le travail lui-même — explorer quatre sources était
nécessaire pour comprendre ce qui avait de la valeur — mais l'\emph{absence de
critère explicite} pendant cette exploration. Il a fallu construire le critère
d'abord, puis mesurer le travail à l'aune de ce critère.

Si c'était à refaire, j'écrirais la question \og l'acheteur pourrait-il obtenir
ceci depuis la source ? \fg{} dès la première semaine, avant d'écrire la
première ligne de connecteur.

\subsection{L'instabilité du déploiement}

Le déploiement a consommé une part disproportionnée du temps. Les messages
d'erreur de la plateforme désignaient systématiquement le symptôme, jamais la
cause. Le message \og le motif ne correspond à aucune fonction \fg{} a coûté
plusieurs itérations alors que la cause réelle — la visibilité du nom racine
dans le module — n'était mentionnée nulle part.

La leçon : \textbf{sur une plateforme d'hébergement, il faut diagnostiquer par
l'expérience contrôlée, pas par le message d'erreur}. Réduire la surface,
changer une variable à la fois, observer.

\subsection{Le périmètre du mode VFX}

Le mode VFX a été spécifié, encadré, partiellement implémenté — et non terminé.
Trois tâches sur dix restent ouvertes. C'est un périmètre trop ambitieux pour
la durée disponible, et la conséquence est un demi-livrable.

Le point positif est que la partie livrée est la partie \emph{gouvernance} :
l'amendement, l'étiquetage, le bus d'entrée et les garanties. C'est la partie
qui protège le produit. Le rendu 3D lui-même est la partie amusante, et c'est
celle qui manque.

\section{Ce que je ferais différemment}

\begin{table}[H]
\centering
\small
\caption{Choix passés et alternatives envisagées.}
\label{tab:retro}
\begin{tabularx}{\textwidth}{@{}X X@{}}
\toprule
\textbf{Ce qui a été fait} & \textbf{Ce que je ferais autrement} \\
\midrule
Explorer quatre sources avant de définir le critère de valeur &
  Écrire le critère de valeur dès le premier jour, puis explorer \\
\addlinespace
Écrire les tests de scénarios après la mise en service &
  Écrire un test de valeur attendue (clé $\rightarrow$ pic $\rightarrow$
  profondeur) dès l'implémentation du moteur \\
\addlinespace
Traiter le déploiement au fil de l'eau &
  Déployer une version minimale dès la première semaine, pour découvrir
  l'infrastructure tôt plutôt que tard \\
\addlinespace
Spécifier dix tâches VFX &
  Spécifier trois tâches VFX (bascule, étiquetage, garanties) et garder le
  rendu pour une phase ultérieure \\
\addlinespace
Garder la calibration de latence pour la fin &
  Mesurer dès que le pipeline est stable, pour que le budget soit un
  engagement tenu et non une hypothèse \\
\bottomrule
\end{tabularx}
\end{table}

\section{Compétences acquises}

\begin{table}[H]
\centering
\small
\caption{Compétences mobilisées et développées.}
\label{tab:competences}
\begin{tabularx}{\textwidth}{@{}l X@{}}
\toprule
\textbf{Domaine} & \textbf{Acquis} \\
\midrule
Systèmes d'information géographique &
  WFS et GML 3.2, test point-in-polygon par lancer de rayons, proximité
  géométrique, projections et ordre des axes, boîtes englobantes \\
\addlinespace
Ingénierie backend &
  FastAPI, Pydantic v2, asynchrone avec \code{httpx}, conception de contrat de
  données, gestion de cache et de quotas \\
\addlinespace
Ingénierie frontend &
  React et TypeScript, Mapbox GL JS (extrusion, couches personnalisées,
  terrain), animation synchronisée à un état applicatif \\
\addlinespace
Qualité logicielle &
  Tests hors-ligne avec fixtures, tests de dégradation, tests de propriété,
  portes de recherche CI \\
\addlinespace
Méthodologie &
  Développement piloté par les spécifications, gouvernance par document,
  collaboration structurée avec des agents de programmation \\
\addlinespace
Déploiement &
  Architecture serverless mono-domaine, réécritures et proxies, diagnostic
  d'incidents de plateforme \\
\addlinespace
Domaine métier &
  Souscription, Solvabilité II, régime CatNat, non-stationnarité du risque
  climatique, distinction entre donnée et jugement \\
\bottomrule
\end{tabularx}
\end{table}

\section{Appréciation personnelle}

Ce stage m'a surtout appris qu'un projet technique se joue moins sur la
difficulté du code que sur la justesse des décisions. Écrire un parseur GML
est un exercice ; décider qu'il ne faut \emph{pas} produire de score est une
décision. La seconde a demandé plus de courage et a eu plus d'impact.

J'en retiens une posture : chercher la formulation la plus honnête possible du
problème, et s'y tenir même quand elle impose de supprimer du travail. Le
recentrage du 25 août a été inconfortable sur le moment — il fallait admettre
que plusieurs semaines de production ne servaient pas le produit — mais c'est
ce qui a donné au livrable sa cohérence.

Le projet n'est pas terminé. Il est, en revanche, défendable : chaque choix est
documenté, chaque fait porte sa source, et chaque limite connue est écrite
plutôt que cachée. C'est probablement la définition d'un travail honnête.

% =====================================================================
\chapter{Conclusion et perspectives}
\label{ch:conclusion}
% =====================================================================

\section{Rappel de la problématique}

Ce stage partait d'un constat simple : les données publiques de risque
climatique existent, mais la jointure entre l'exposition d'un bâtiment et sa
vulnérabilité n'existe nulle part. Un souscripteur qui doit évaluer un dossier
le fait aujourd'hui à l'œil, en croisant deux portails, sans traçabilité.

La question posée était donc : peut-on construire, à partir de sources
publiques uniquement, un service qui réponde à \og ce bâtiment est-il exposé, et
qu'est-ce qui le rend vulnérable ? \fg{}, avec une qualité d'information
suffisante pour qu'un professionnel de l'assurance s'appuie dessus ?

\section{Réponse apportée}

La réponse est oui, avec des conditions. Le stage a produit un service
fonctionnel, déployé et vérifié, qui :

\begin{enumerate}[leftmargin=1.6em]
  \item résout une adresse française et interroge les sources en parallèle ;
  \item teste le point résolu contre les périmètres d'aléas par géométrie
        réelle, pas par approximation communale ;
  \item associe à chaque aléa la description verbatim du bâtiment issue de la
        BDNB ;
  \item porte, sur chaque fait, sa source, sa date et son niveau de résolution ;
  \item refuse explicitement de produire un score ou une recommandation ;
  \item distingue trois états — présent, absent, inconnu — de sorte qu'une
        panne ne se lise jamais comme une absence de risque ;
  \item met le résultat en situation sur une carte, avec un moteur de scénarios
        fondé sur les classes réglementaires, sans jamais inventer un niveau ;
  \item est couvert par 357 tests automatisés hors-ligne.
\end{enumerate}

Les conditions sont aussi importantes que la réponse. Le service n'est pas
commercialisable en l'état : l'authentification est simulée, le magasin de lots
est en mémoire, la certification attend sa qualification humaine, et le budget
de latence n'est pas encore calibré. Le prototype démontre une thèse ; il ne
remplace pas un produit.

\section{Les trois idées à retenir}

\subsection{Le refus du score est une position tenable}

L'intuition initiale — produire un score global de risque — a été délibérément
abandonnée. Ce n'est pas une concession, c'est un positionnement : sous
Solvabilité~II, un score tiers engage son producteur, se fait reconstruire ou
refuser par l'utilisateur, et mélange des aléas qui ne se compensent pas.

Le produit a été reconstruit sur d'autres fondations — provenance et résolution
— et s'est révélé démontrable sans son chiffre. Le champ de résolution, qui
n'était qu'un attribut parmi d'autres, est devenu le signal de qualité central.

\subsection{L'honnêteté sur l'incertitude est une fonctionnalité}

La trichotomie \code{true}/\code{false}/\code{null}, la distinction entre
absence et panne, l'étiquetage du mode cinématique : tout cela relève du même
principe. Un produit de données ne se juge pas seulement sur ce qu'il affirme,
mais sur la manière dont il dit \og je ne sais pas \fg.

C'est aussi ce qui rend le produit utilisable en contexte réglementaire : une
décision peut être justifiée si l'on sait quel fait a été établi à quel niveau.

\subsection{La soustraction est un travail d'ingénierie}

La phase la plus productive du stage a été celle où le périmètre a été réduit.
Supprimer quatre connecteurs, une couche de notation, un tableau de bord et des
milliers de lignes a augmenté la valeur du livrable. La règle \og pas de code
dormant \fg{} et l'historique Git comme unique archive ont rendu cette
soustraction possible sans perte d'information.

\section{Perspectives}

Le travail identifié pour la suite se répartit en trois horizons.

\subsection{Court terme — finir et qualifier}

\begin{itemize}[leftmargin=1.4em]
  \item \textbf{Qualifier la certification.} Confronter les fixtures des
        quatre aléas certifiés aux cartes officielles et signer la version~1.
  \item \textbf{Calibrer la latence.} Exécuter la campagne de mesures sur au
        moins 100 appels réels et décider, sur données, si la parallélisation
        des couches WFS est nécessaire.
  \item \textbf{Exécuter le test manuel de bout en bout.} Diagnostiquer une
        adresse certifiée TRI un jour de pluie et valider la boucle complète.
  \item \textbf{Terminer le mode VFX} — ou acter qu'il reste un démonstrateur
        partiel, ce qui est défendable puisque la partie gouvernance est livrée.
\end{itemize}

\subsection{Moyen terme — rendre le service exploitable}

\begin{itemize}[leftmargin=1.4em]
  \item \textbf{Authentification réelle} — remplacer le mode démonstration par
        un système de comptes, condition de tout usage externe.
  \item \textbf{Persistance du store de lot} — remplacer le magasin en mémoire
        par une file de traitement, seule façon de tenir le volume.
  \item \textbf{Cache partagé} — sortir le cache du processus pour qu'il
        survive au serverless.
  \item \textbf{Source vectorielle pour le RGA} — câbler une source
        polygonale pour le retrait-gonflement des argiles, ce qui exige un
        amendement constitutionnel, et supprimerait la dernière étiquette
        \code{commune-level-estimate} structurelle.
\end{itemize}

\subsection{Long terme — le produit complet}

\begin{itemize}[leftmargin=1.4em]
  \item \textbf{La trajectoire climatique.} C'est l'argument de fond du
        produit : la non-stationnarité du risque. Le prototype en pose la
        structure — horizons explicites, scénarios étiquetés, unités réelles —
        mais il reste à industrialiser la lecture des projections et à exposer
        les variables par aléa et par horizon dans le contrat.
  \item \textbf{La livraison en lot pour l'actuaire.} Le second acheteur
        attend de la donnée structurée à l'échelle d'un portefeuille, avec les
        mêmes sémantiques que le produit interactif. L'invariant \og le lot est
        N fois l'unitaire \fg{} a été posé précisément pour rendre cette
        extension non coûteuse.
  \item \textbf{La certification étendue.} Passer de quatre aléas certifiés à
        l'ensemble des aléas couverts par un périmètre vectoriel.
\end{itemize}

\section{Mot de la fin}

\typhoon{} est un prototype, et il faut le présenter comme tel. Mais c'est un
prototype qui démontre quelque chose de précis : \emph{la donnée publique
française permet de construire une information de risque climatique par
bâtiment, traçable et défendable} — à condition d'accepter de ne pas tout
savoir et de le dire.

Le stage m'a laissé une conviction : dans un domaine où la donnée est
abondante, la valeur ne se crée pas en ajoutant des sources ou des indicateurs,
mais en établissant avec rigueur une jointure que personne n'avait faite, et en
assumant explicitement les limites de ce que l'on affirme.

\begin{quote}
\itshape
Un produit de données honnête ne dit pas seulement ce qu'il sait.\\
Il dit aussi comment il le sait, et à quel point il le sait.
\end{quote}


% ---------------------------------------------------------------------
%  Annexes
% ---------------------------------------------------------------------
\appendix
% =====================================================================
\chapter{Référence du contrat canonique}
\label{ann:contrat}
% =====================================================================

Cette annexe reprend la spécification complète du contrat de données, telle
qu'elle figure dans la spécification du projet. Elle est destinée à un
intégrateur qui souhaite consommer l'API.

\section{Structure complète}

\begin{lstlisting}[style=json,caption={Contrat canonique \texttt{DiagnosticRecord}
v1.0.},label={lst:contrat-full}]
POST /diagnostic/adresse
body: { "adresse": "10 rue de Rivoli, 75004 Paris" }

200 ->
{
  "schema_version": "1.0",
  "adresse": {
    "saisie": "...",              // entree brute
    "normalisee": "...",          // libelle du geocodeur
    "citycode": "75056",
    "postcode": "75004",
    "city": "Paris",
    "lat": 48.8566,
    "lon": 2.3522,
    "geocode_score": 0.95
  },
  "aleas": [                      // 13 objets, referentiel complet
    {
      "code": "inondation",
      "libelle": "Inondation",
      "present": true,            // true | false | null
      "present_commune": true,
      "zonage": "Dans un perimetre PPR inondation",   // texte verbatim
      "hauteur_eau_m": 1.5,       // optionnels specifiques a l'alea,
      "zone_sismique": null,      // TOUJOURS presents, null si non applicable
      "catnat_historique": [],
      "per_building": {           // present uniquement si teste geometriquement
        "method": "point-in-polygon",  // "point-in-polygon" | "proximity"
        "radius_m": null,              // echo de la proximite si applicable
        "count": 2
      },
      "source": "georisques",
      "url_detail": "https://www.georisques.gouv.fr/risques/inondations",
      "erreur": null,             // string | null
      "resolution": "per-building"
      // "per-building" | "commune-level" | "commune-level-estimate"
    }
  ],
  "bdnb": {                       // bloc entier null si BDNB indisponible
    "donnees": {
      // ~139 champs BDNB verbatim
      // batiment = batiment principal
      // autres_batiments_meme_adresse = [...]
    },
    "_source": {
      "provider": "BDNB",
      "url": "https://api.bdnb.io",
      "attribution": "Source BDNB (CSTB) - donnees a jour au <millesime>",
      "recuperee_le": "2026-08-25T10:00:00+00:00"
    }
  },
  "georisques_source": {
    "provider": "Georisques",
    "url": "https://www.georisques.gouv.fr",
    "attribution": "Source Georisques (BRGM / MTE) - donnees a jour au <millesime>",
    "recuperee_le": "2026-08-25T10:00:00+00:00"
  },
  "erreurs_partielles": [],
  "genere_le": "2026-08-25T10:00:01+00:00",
  "avertissement": "Ce rapport agrege les donnees publiques Georisques..."
}
\end{lstlisting}

\section{Champs interdits}

Les champs suivants ne peuvent pas exister dans le contrat. Leur présence fait
échouer la validation du schéma, et une porte de recherche automatique échoue si
l'un d'eux apparaît dans un schéma émis :

\begin{itemize}[leftmargin=1.4em]
  \item \code{niveau} et les bandes D03 ;
  \item tout score numérique, agrégé ou composite ;
  \item \code{recommandations} et toute évaluation générée par IA ;
  \item \code{copernicus} dans le contrat canonique v1.
\end{itemize}

\section{Règles de sérialisation}

\begin{table}[H]
\centering
\small
\caption{Règles de sérialisation à respecter par tout consommateur.}
\label{tab:serialisation}
\begin{tabularx}{\textwidth}{@{}l X@{}}
\toprule
\textbf{Règle} & \textbf{Énoncé} \\
\midrule
Optionnels toujours présents &
  Les champs spécifiques à un aléa sont toujours présents, à \code{null}
  lorsqu'ils ne s'appliquent pas. Une clé n'est jamais absente. \\
\addlinespace
Exception \code{per\_building} &
  L'objet \code{per\_building} n'existe que si le point a été testé
  géométriquement. Sa présence porte l'information. \\
\addlinespace
BDNB verbatim &
  Les clés BDNB reçues et émises sont identiques — vérifié par test de
  conformité. \\
\addlinespace
Attribution non vide &
  Chaque bloc \code{\_source} porte une chaîne d'attribution non vide. \\
\addlinespace
Résolution contrainte &
  Le RGA ne peut sérialiser aucune autre valeur que
  \code{commune-level-estimate}. \\
\bottomrule
\end{tabularx}
\end{table}

\section{Le référentiel des treize aléas}

\begin{table}[H]
\centering
\footnotesize
\caption{Les treize aléas, leur libellé et leur niveau de résolution attendu.}
\label{tab:aleas}
\begin{tabularx}{\textwidth}{@{}l l X@{}}
\toprule
\textbf{Code} & \textbf{Libellé} & \textbf{Résolution attendue} \\
\midrule
\code{icpe} & Installations classées &
  \code{per-building} ou \code{commune-level} \\
\code{inondation} & Inondation & \code{per-building} \\
\code{sismicite} & Sismicité & \code{commune-level} (décret, permanent) \\
\code{mouvement\_terrain} & Mouvement de terrain & \code{per-building} \\
\code{radon} & Radon & \code{commune-level} (permanent) \\
\code{rga} & Retrait-gonflement des argiles &
  \code{commune-level-estimate} (forcé) \\
\code{cavite} & Cavités souterraines & \code{per-building} ou communal \\
\code{feu\_foret} & Feu de forêt & \code{per-building} \\
\code{avalanche} & Avalanche & \code{per-building} \\
\code{canalisations} & Canalisations (TMD) &
  \code{per-building} par proximité (200~m) \\
\code{vent\_cyclonique} & Vent cyclonique & \code{commune-level} \\
\code{ppr} & PPR (agrégat) & \code{per-building} \\
\code{ssp} & Sites et sols pollués &
  \code{per-building} par proximité (200~m) \\
\bottomrule
\end{tabularx}
\end{table}

% =====================================================================
\chapter{Surface de l'API et mise en œuvre}
\label{ann:api}
% =====================================================================

\section{Points d'entrée}

\begin{table}[H]
\centering
\footnotesize
\caption{Points d'entrée de l'API, avec authentification requise.}
\label{tab:api-full}
\begin{tabularx}{\textwidth}{@{}l l X@{}}
\toprule
\textbf{Méthode et chemin} & \textbf{Auth} & \textbf{Description} \\
\midrule
\api{POST /diagnostic/adresse} & oui &
  Transaction canonique, adresse $\rightarrow$ enregistrement complet \\
\api{POST /diagnostic/batch} & oui &
  Soumission d'un lot ($\leq$ 10\,000 adresses) \\
\api{GET /diagnostic/batch/\{id\}} & oui &
  État et résultats du lot \\
\api{GET /diagnostic/adresse/rapport-pdf} & oui &
  Relais vers le PDF officiel Géorisques \\
\api{GET /diagnostic/zone/building} & oui &
  Fiche BDNB au point \\
\api{GET /diagnostic/zone/buildings} & oui &
  Fiches BDNB autour du point \\
\api{GET /api/flood-alea} & oui &
  Classes TRI au point \\
\api{GET /api/hydro} & oui & Tracé hydrographique \\
\api{GET /api/hydro/extend} & oui & Extension du tracé \\
\api{GET /api/meteo} & oui & Pluie horaire et débit \\
\api{GET /api/photo} & oui & Photographie Panoramax et crédit \\
\api{GET /api/report} & oui & Rapport narratif \\
\api{GET /api/report/stream} & oui & Rapport narratif en flux \\
\api{GET /api/geocode/search} & oui & Autocomplétion d'adresse \\
\api{GET /health} & non & Sonde de vie \\
\api{GET /health/detailed} & non & Sonde détaillée \\
\bottomrule
\end{tabularx}
\end{table}

\section{Format des erreurs}

\begin{table}[H]
\centering
\small
\caption{Codes d'erreur et signification.}
\label{tab:erreurs-api}
\begin{tabularx}{\textwidth}{@{}l l X@{}}
\toprule
\textbf{Code} & \textbf{Identifiant} & \textbf{Signification} \\
\midrule
401 & — & Clé d'API absente \\
403 & — & Clé d'API invalide \\
422 & \code{adresse\_non\_trouvee} & Adresse introuvable par le géocodeur \\
422 & \code{adresse\_ambigue} & Score de géocodage inférieur à 0,4 \\
429 & \code{budget\_epuise} & Budget mensuel amont épuisé (lot uniquement) \\
502 & \code{geocodage\_indisponible} & Géocodeur injoignable ou saturé \\
502 & — & Exception inattendue (jamais de 500 HTML) \\
\bottomrule
\end{tabularx}
\end{table}

\section{Commandes d'installation}

\begin{lstlisting}[language=bash,caption={Installation et exécution en local.},
label={lst:install}]
# Backend (Python 3.11+)
cd backend
python -m venv venv && source venv/bin/activate
pip install -r requirements.txt
uvicorn app.main:app --reload --port 8000

# Frontend (Node 18+)
cd frontend
npm install
npm run dev        # http://localhost:5173
\end{lstlisting}

Le fichier d'exemple d'environnement, à copier à la racine du dépôt,
documente chaque variable nécessaire : jeton Mapbox, URL et clé Supabase, clé
Windy, clé du modèle de rédaction, identifiants Copernicus, et clés d'API
partenaires.

\section{Commandes de vérification}

\begin{lstlisting}[language=bash,caption={Exécution des suites de test.},
label={lst:tests-cmd}]
# Backend (depuis backend/)
python -m pytest tests/ -q     # 214 tests hors-ligne, < 60 s
ruff check app tests           # lint

# Frontend (depuis frontend/)
npm test                       # 143 tests Vitest
npm run build                  # tsc -b && vite build
\end{lstlisting}

\section{Vérification de l'instance déployée}

\begin{lstlisting}[language=bash,caption={Contrôles de bon fonctionnement de
l'instance en ligne.},label={lst:verif-cmd}]
BASE=https://<projet>.vercel.app
curl $BASE/health
curl "$BASE/api/flood-alea?lat=48.848&lon=2.370" | head -c 200
curl -X POST "$BASE/diagnostic/adresse" \
  -H 'Content-Type: application/json' \
  -d '{"adresse":"Quai de la Rapee, 75012 Paris"}' | head -c 200
\end{lstlisting}

Le réveil d'une fonction froide prend de 5 à 10~s, et un diagnostic complet
peut atteindre environ 15~s à froid. C'est normal sur l'offre gratuite :
appeler la sonde de vie une fois avant une démonstration suffit à préchauffer
l'instance.

\section{Reproduction des captures d'écran}

\begin{lstlisting}[language=bash,caption={Génération des captures
d'écran.},label={lst:shots}]
cd frontend && node scripts/demo-shots.mjs
# necessite le backend et le serveur vite en local
\end{lstlisting}

% =====================================================================
\chapter{Glossaire}
\label{ann:glossaire}
% =====================================================================

\section{Termes métier}

\begin{description}[leftmargin=0em,style=nextline]
  \item[Souscription (\emph{underwriting}).] Processus par lequel un assureur
  accepte, met à l'étude ou refuse de couvrir un risque, et détermine les
  conditions de la couverture.
  \item[Souscripteur.] Professionnel qui traite les dossiers de souscription au
  quotidien.
  \item[Actuaire.] Professionnel qui construit les modèles de tarification et
  de provisionnement, en s'appuyant sur des données plutôt que sur une
  interface.
  \item[Solvabilité II.] Cadre réglementaire européen de solvabilité des
  assureurs. Il impose notamment la traçabilité et la justification des
  décisions.
  \item[Régime CatNat.] Régime français d'indemnisation des catastrophes
  naturelles, fondé sur la reconnaissance de l'état de catastrophe naturelle
  par arrêté interministériel.
  \item[Non-stationnarité.] Propriété d'un risque dont la distribution change
  dans le temps. Le risque climatique est non stationnaire : le passé ne prédit
  pas directement le futur.
  \item[PPR.] Plan de Prévention des Risques. Document réglementaire qui
  délimite des zones exposées à un aléa et fixe des prescriptions.
  \item[TRI.] Territoire à Risque Important d'inondation, issu de la Directive
  Inondation. Classe l'exposition en quatre niveaux de période de retour.
  \item[Période de retour.] Durée moyenne entre deux occurrences d'un événement
  d'intensité donnée. Une crue centennale a une probabilité d'environ 1\,\% par
  an.
  \item[RGA.] Retrait-gonflement des argiles. Phénomène de déformation des sols
  argileux selon leur humidité, cause majeure de fissuration des bâtiments.
  \item[CatNat historique.] Liste des arrêtés de catastrophe naturelle pris
  pour une commune.
\end{description}

\section{Termes techniques}

\begin{description}[leftmargin=0em,style=nextline]
  \item[WFS.] Web Feature Service. Standard de l'Open Geospatial Consortium
  pour l'échange de données géographiques vectorielles.
  \item[GML.] Geography Markup Language. Format XML d'encodage de géométries,
  utilisé par les services WFS. Le projet traite GML 3.2.
  \item[Point-in-polygon.] Test déterminant si un point est à l'intérieur d'un
  polygone. Implémenté par lancer de rayon.
  \item[Lancer de rayon.] Algorithme consistant à compter les intersections
  d'un rayon partant du point avec les arêtes du polygone. Un nombre impair
  signifie que le point est à l'intérieur.
  \item[Boîte englobante (BBOX).] Rectangle délimitant la zone d'une requête
  spatiale.
  \item[Proximité géométrique.] Test applicable aux géométries non polygonales
  (points, lignes) : le point est-il à moins d'un rayon donné ?
  \item[Ordre des axes.] Convention de notation des coordonnées
  (latitude/longitude ou longitude/latitude). Une erreur d'ordre fait échouer
  silencieusement tout test géométrique.
  \item[Verbatim.] Passage d'une donnée sans transformation : ni renommage, ni
  filtrage, ni recalcul.
  \item[Provenance.] Métadonnée décrivant l'origine d'un fait : fournisseur,
  URL, chaîne d'attribution, horodatage de collecte.
  \item[Résolution.] Niveau auquel un fait a été établi : par bâtiment, par
  commune, ou estimé par commune.
  \item[Mode dégradé.] Comportement du service lorsqu'une source échoue : la
  réponse reste valide mais partielle, avec les faits concernés marqués comme
  inconnus.
  \item[Fixture.] Donnée de test figée, capturée une fois depuis une source
  réelle, permettant des tests hors-ligne reproductibles.
  \item[Serverless.] Modèle d'exécution où le code s'exécute dans une fonction
  éphémère, sans serveur permanent. Le système de fichiers n'y persiste pas.
  \item[Étiquetage.] Mention visible signalant qu'un rendu est illustratif et
  non contractuel.
  \item[Porte de recherche (\emph{grep gate}).] Vérification automatique qui
  échoue si un motif interdit (par exemple un champ de score) apparaît dans le
  code ou un schéma émis.
\end{description}

\section{Sigles}

\begin{table}[H]
\centering
\small
\caption{Sigles et acronymes utilisés dans le rapport.}
\label{tab:sigles}
\begin{tabularx}{\textwidth}{@{}l X@{}}
\toprule
\textbf{Sigle} & \textbf{Signification} \\
\midrule
API & Application Programming Interface \\
BAN & Base Adresse Nationale \\
BDNB & Base de Données Nationale des Bâtiments \\
BBOX & \emph{Bounding box} — boîte englobante \\
CDS & Climate Data Store (Copernicus) \\
CORS & \emph{Cross-Origin Resource Sharing} \\
CSTB & Centre Scientifique et Technique du Bâtiment \\
DPE & Diagnostic de Performance Énergétique \\
GML & Geography Markup Language \\
IGN & Institut Géographique National \\
INSEE & Institut National de la Statistique et des Études Économiques \\
LLM & \emph{Large Language Model} — grand modèle de langage \\
OGC & Open Geospatial Consortium \\
PPR & Plan de Prévention des Risques \\
PPRI & Plan de Prévention du Risque Inondation \\
RCP & \emph{Representative Concentration Pathway} — trajectoire de
  concentration représentative (scénario climatique) \\
REST & \emph{Representational State Transfer} \\
RGA & Retrait-gonflement des argiles \\
SSP & Sites et sols pollués \\
SPA & \emph{Single Page Application} \\
TMD & Transport de matières dangereuses \\
TRI & Territoire à Risque Important d'inondation \\
TTL & \emph{Time To Live} — durée de vie (d'un cache) \\
VFX & \emph{Visual effects} — effets visuels (mode de rendu cinématique) \\
WFS & Web Feature Service \\
WMS & Web Map Service \\
\bottomrule
\end{tabularx}
\end{table}


% ---------------------------------------------------------------------
%  Bibliographie
% ---------------------------------------------------------------------
% =====================================================================
\chapter*{Sources et références}
\addcontentsline{toc}{chapter}{Sources et références}
\markboth{Sources et références}{Sources et références}
% =====================================================================

\section*{Sources de données utilisées}

\begin{itemize}[leftmargin=1.4em]
  \item \textbf{Géorisques} — portail national des risques naturels et
        technologiques, Ministère de la Transition écologique / BRGM.
        Données distribuées sous Licence Ouverte 2.0 (Etalab).
        \url{https://www.georisques.gouv.fr}
  \item \textbf{BDNB} — Base de Données Nationale des Bâtiments, Centre
        Scientifique et Technique du Bâtiment (CSTB). Licence Ouverte 2.0.
        \url{https://api.bdnb.io}
  \item \textbf{Géoplateforme} — service de géocodage, Institut Géographique
        National (IGN). \url{https://data.geopf.fr/geocodage}
  \item \textbf{Sandre} — portail national d'accès aux données sur l'eau.
        \url{https://www.sandre.eaufrance.fr}
  \item \textbf{Open-Meteo} — API météorologique ouverte, incluant le modèle
        hydrologique GloFAS. \url{https://open-meteo.com}
  \item \textbf{Panoramax} — service public de photographie de terrain
        contributive. \url{https://panoramax.fr}
  \item \textbf{Copernicus Climate Data Store} — projections climatiques
        européennes, ECMWF. Jeu de données \emph{Climate indicators for Europe
        from 1940 to 2100}. \url{https://cds.climate.copernicus.eu}
  \item \textbf{Vigicrues} — service public de vigilance crues.
        \url{https://www.vigicrues.gouv.fr}
  \item \textbf{API Carto / BAN} — Base Adresse Nationale.
        \url{https://adresse.data.gouv.fr}
\end{itemize}

\section*{Cadres réglementaires}

\begin{itemize}[leftmargin=1.4em]
  \item \textbf{Licence Ouverte 2.0} (Etalab) — cadre juridique de
        réutilisation des données publiques françaises. Réutilisation
        commerciale autorisée, obligation d'attribution.
        \url{https://www.etalab.gouv.fr/licence-ouverte-open-licence}
  \item \textbf{Directive Inondation} — cadre européen ayant conduit à la
        définition des Territoires à Risque Important d'inondation (TRI).
  \item \textbf{Régime CatNat} — dispositif français d'indemnisation des
        catastrophes naturelles, article L125-1 du Code des assurances.
  \item \textbf{Solvabilité II} — directive européenne 2009/138/CE encadrant la
        solvabilité des organismes d'assurance et de réassurance.
\end{itemize}

\section*{Standards techniques}

\begin{itemize}[leftmargin=1.4em]
  \item \textbf{OGC Web Feature Service 2.0} — standard d'échange de données
        géographiques vectorielles.
  \item \textbf{OGC Geography Markup Language 3.2} — format d'encodage des
        géométries.
  \item \textbf{Mapbox GL JS} — bibliothèque de rendu cartographique.
  \item \textbf{FastAPI} — cadre de développement d'API Python.
  \item \textbf{Pydantic v2} — validation de données Python.
\end{itemize}

\section*{Documents internes au projet}

\begin{itemize}[leftmargin=1.4em]
  \item \fichier{constitution.md} — règles non négociables, version 0.2,
        amendée le 14 septembre 2026.
  \item \fichier{docs/workflow/01-brainstorm-notes.md} — journal des décisions
        de la phase 1.
  \item \fichier{docs/workflow/concept.md} — concept et questions ouvertes
        résolues (OQ-1 à OQ-6).
  \item \fichier{docs/workflow/spec.md} — spécification fonctionnelle,
        28 exigences \frref{01} à \frref{28}.
  \item \fichier{docs/workflow/plan.md} — plan d'exécution, tâches T001 à T018.
  \item \fichier{docs/workflow/scenario-plan.md} — plan du moteur de scénarios,
        tâches \scn{001} à \scn{050}.
  \item \fichier{docs/workflow/vfx-spec.md} et \fichier{vfx-plan.md} —
        spécification et plan du mode de rendu cinématique.
  \item \fichier{AGENTS.md} — conventions, invariants acquis et pièges
        techniques.
  \item \fichier{handoff.md} — décisions architecturales en attente.
  \item \fichier{roadmapproduction.md} — feuille de route vers la production.
  \item \fichier{README.md} — documentation d'installation, d'usage et de
        déploiement.
\end{itemize}


\end{document}
